% Generated from data/spectral-lines.json. Do not edit.
\section{H I 21 cm}
\index{H I}

\begin{conceptbox}[At a glance]
\begin{tabularx}{\textwidth}{@{}lX@{}}
\textbf{Rest coordinate} & 1420.405751768 MHz \\
\textbf{Band} & radio / L band \\
\textbf{Transition} & 1s ground-state hyperfine F=1 -> 0 \\
\textbf{Mechanism} & magnetic-dipole hyperfine transition \\
\textbf{Primary tracers} & neutral atomic hydrogen column; gas kinematics; rotation curves; spin temperature with absorption; large-scale structure via intensity mapping \\
\textbf{Typical environments} & cold and warm neutral medium; galaxy disks; high-velocity clouds; damped Ly-alpha systems in absorption; reionization-era IGM statistically \\
\textbf{Indicative regime} & CNM roughly 20-200 K; WNM roughly 5000-10000 K; detectable outside strict LTE \\
\textbf{Density guidance} & Not a conventional density diagnostic. Emission gives optically thin H I column; emission plus absorption constrains spin temperature and covering fraction. \\
\end{tabularx}
\end{conceptbox}

H I 21 cm is produced by magnetic-dipole hyperfine transition and is encountered in cold and warm neutral medium; galaxy disks; high-velocity clouds; damped Ly-alpha systems in absorption; reionization-era IGM statistically. It is commonly used to trace neutral atomic hydrogen column; gas kinematics; rotation curves; spin temperature with absorption; large-scale structure via intensity mapping.

\subsection*{How to use it}
Use this transition with its diagnostic companions to constrain neutral atomic hydrogen column; gas kinematics; rotation curves; spin temperature with absorption; large-scale structure via intensity mapping, while fitting the continuum, line-spread function, and relevant blends.

\subsection*{Interpretive cautions}
Rest frequency is exceptionally precise. State the radio/optical/relativistic velocity convention and reference frame. Self-absorption and optical depth can hide cold gas.

\subsection*{Useful companions}
Ly-alpha 1215.6701 A; CO(1-0) 115.2712018 GHz; OH 1665/1667 MHz

\subsection*{Sources}
\texttt{NRAO\_HI\_OH; NIST\_ASD; DRAINE\_2011}

\clearpage


\section{OH 1665 MHz main line}
\index{OH}

\begin{conceptbox}[At a glance]
\begin{tabularx}{\textwidth}{@{}lX@{}}
\textbf{Rest coordinate} & 1665.4018 MHz \\
\textbf{Band} & radio / 18 cm \\
\textbf{Transition} & ground-state 2Pi3/2, J=3/2 hyperfine main-line transition \\
\textbf{Mechanism} & lambda-doubling plus hyperfine transition; often masing \\
\textbf{Primary tracers} & molecular gas magnetic field via Zeeman splitting; maser kinematics; pumped dense gas \\
\textbf{Typical environments} & massive star-forming regions; circumstellar envelopes; compact H II regions; megamaser hosts \\
\textbf{Indicative regime} & typically tens to hundreds of kelvin; non-LTE pumping controls masers \\
\textbf{Density guidance} & Population inversion depends on FIR pump, OH column and density; not a simple one-number density tracer. \\
\end{tabularx}
\end{conceptbox}

OH 1665 MHz main line is produced by lambda-doubling plus hyperfine transition; often masing and is encountered in massive star-forming regions; circumstellar envelopes; compact H II regions; megamaser hosts. It is commonly used to trace molecular gas magnetic field via Zeeman splitting; maser kinematics; pumped dense gas.

\subsection*{How to use it}
Use this transition with its diagnostic companions to constrain molecular gas magnetic field via Zeeman splitting; maser kinematics; pumped dense gas, while fitting the continuum, line-spread function, and relevant blends.

\subsection*{Interpretive cautions}
Compare 1665:1667 and satellite-line ratios to diagnose LTE departures, pumping and opacity. Strongly affected by radio-frequency interference in some bands.

\subsection*{Useful companions}
OH 1667 MHz; OH satellite lines 1612 and 1720 MHz; H2O 22 GHz maser

\subsection*{Sources}
\texttt{NRAO\_HI\_OH; JPL; SPLATALOGUE}

\clearpage


\section{OH 1667 MHz main line}
\index{OH}

\begin{conceptbox}[At a glance]
\begin{tabularx}{\textwidth}{@{}lX@{}}
\textbf{Rest coordinate} & 1667.35903 MHz \\
\textbf{Band} & radio / 18 cm \\
\textbf{Transition} & ground-state 2Pi3/2, J=3/2 hyperfine main-line transition \\
\textbf{Mechanism} & lambda-doubling plus hyperfine transition; thermal absorption/emission or maser \\
\textbf{Primary tracers} & molecular gas; dust-obscured mergers; magnetic fields; line-of-sight velocity \\
\textbf{Typical environments} & diffuse molecular clouds; star-forming regions; OH megamaser galaxies; circumstellar envelopes \\
\textbf{Indicative regime} & typically 10-300 K, with excitation set by collisions and radiative pumping \\
\textbf{Density guidance} & Line ratios are excitation/opacity diagnostics rather than direct density meters. \\
\end{tabularx}
\end{conceptbox}

OH 1667 MHz main line is produced by lambda-doubling plus hyperfine transition; thermal absorption/emission or maser and is encountered in diffuse molecular clouds; star-forming regions; OH megamaser galaxies; circumstellar envelopes. It is commonly used to trace molecular gas; dust-obscured mergers; magnetic fields; line-of-sight velocity.

\subsection*{How to use it}
Use this transition with its diagnostic companions to constrain molecular gas; dust-obscured mergers; magnetic fields; line-of-sight velocity, while fitting the continuum, line-spread function, and relevant blends.

\subsection*{Interpretive cautions}
The 1667-MHz line usually dominates extragalactic OH megamasers. Hyperfine LTE expectations can fail spectacularly under FIR pumping.

\subsection*{Useful companions}
OH 1665 MHz; OH 1612/1720 MHz; FIR OH pump transitions

\subsection*{Sources}
\texttt{NRAO\_HI\_OH; JPL; SPLATALOGUE}

\clearpage


\section{H2O 22 GHz maser}
\index{H2O}

\begin{conceptbox}[At a glance]
\begin{tabularx}{\textwidth}{@{}lX@{}}
\textbf{Rest coordinate} & 22.23508 GHz \\
\textbf{Band} & microwave / K band \\
\textbf{Transition} & 6\_16 -> 5\_23, ortho-H2O \\
\textbf{Mechanism} & collisionally pumped rotational maser \\
\textbf{Primary tracers} & sub-au kinematics; geometric distances; black-hole masses; dense shocked gas \\
\textbf{Typical environments} & shocked protostellar outflows; circumstellar envelopes; AGN accretion disks; nuclear winds \\
\textbf{Indicative regime} & commonly about 300-1000 K for strong collisional pumping \\
\textbf{Density guidance} & Maser inversion favors very dense gas, often n(H2) about 10\textasciicircum{}7-10\textasciicircum{}10 cm\textasciicircum{}-3, but geometry and velocity coherence dominate brightness. \\
\end{tabularx}
\end{conceptbox}

H2O 22 GHz maser is produced by collisionally pumped rotational maser and is encountered in shocked protostellar outflows; circumstellar envelopes; AGN accretion disks; nuclear winds. It is commonly used to trace sub-au kinematics; geometric distances; black-hole masses; dense shocked gas.

\subsection*{How to use it}
Use this transition with its diagnostic companions to constrain sub-au kinematics; geometric distances; black-hole masses; dense shocked gas, while fitting the continuum, line-spread function, and relevant blends.

\subsection*{Interpretive cautions}
Brightness temperature is non-thermal and cannot be interpreted as kinetic temperature. Narrow spikes vary and may be highly beamed.

\subsection*{Useful companions}
SiO v=1 J=1-0; OH 1665/1667 MHz; thermal sub-mm H2O lines

\subsection*{Sources}
\texttt{JPL; SPLATALOGUE; NRAO\_HI\_OH}

\clearpage


\section{NH3 (1,1) inversion}
\index{NH3}

\begin{conceptbox}[At a glance]
\begin{tabularx}{\textwidth}{@{}lX@{}}
\textbf{Rest coordinate} & 23.6944955 GHz \\
\textbf{Band} & microwave / K band \\
\textbf{Transition} & (J,K)=(1,1) inversion, resolved hyperfine complex \\
\textbf{Mechanism} & metastable inversion transition \\
\textbf{Primary tracers} & dense gas; kinetic temperature with NH3 (2,2); optical depth via satellites; velocity dispersion \\
\textbf{Typical environments} & cold dense cores; infrared dark clouds; molecular clouds; galactic nuclei \\
\textbf{Indicative regime} & especially effective near 8-40 K, but present in warmer dense gas \\
\textbf{Density guidance} & Effective excitation densities are typically around 10\textasciicircum{}3 cm\textasciicircum{}-3 and model dependent; hyperfine anomalies can affect opacity estimates. \\
\end{tabularx}
\end{conceptbox}

NH3 (1,1) inversion is produced by metastable inversion transition and is encountered in cold dense cores; infrared dark clouds; molecular clouds; galactic nuclei. It is commonly used to trace dense gas; kinetic temperature with NH3 (2,2); optical depth via satellites; velocity dispersion.

\subsection*{How to use it}
Use this transition with its diagnostic companions to constrain dense gas; kinetic temperature with NH3 (2,2); optical depth via satellites; velocity dispersion, while fitting the continuum, line-spread function, and relevant blends.

\subsection*{Interpretive cautions}
Fit the full hyperfine structure. The (2,2)/(1,1) rotational temperature is converted to kinetic temperature with a collisional model.

\subsection*{Useful companions}
NH3 (2,2) 23.7226333 GHz; NH3 (3,3) 23.8701296 GHz; N2H+ (1-0)

\subsection*{Sources}
\texttt{JPL; SPLATALOGUE; SHIRLEY\_2015}

\clearpage


\section{NH3 (2,2) inversion}
\index{NH3}

\begin{conceptbox}[At a glance]
\begin{tabularx}{\textwidth}{@{}lX@{}}
\textbf{Rest coordinate} & 23.7226333 GHz \\
\textbf{Band} & microwave / K band \\
\textbf{Transition} & (J,K)=(2,2) inversion, resolved hyperfine complex \\
\textbf{Mechanism} & metastable inversion transition \\
\textbf{Primary tracers} & gas temperature with NH3 (1,1); warmer dense component; kinematics \\
\textbf{Typical environments} & dense molecular cores; massive star-forming clumps; galactic centers \\
\textbf{Indicative regime} & roughly 15-100 K for common applications; higher metastable lines extend the thermometer \\
\textbf{Density guidance} & Similar effective-density scale to (1,1), with excitation dependent on temperature and NH3 column. \\
\end{tabularx}
\end{conceptbox}

NH3 (2,2) inversion is produced by metastable inversion transition and is encountered in dense molecular cores; massive star-forming clumps; galactic centers. It is commonly used to trace gas temperature with NH3 (1,1); warmer dense component; kinematics.

\subsection*{How to use it}
Use this transition with its diagnostic companions to constrain gas temperature with NH3 (1,1); warmer dense component; kinematics, while fitting the continuum, line-spread function, and relevant blends.

\subsection*{Interpretive cautions}
Use matched beams and simultaneous calibration when deriving temperatures; (2,2) may miss the coldest component at low signal-to-noise.

\subsection*{Useful companions}
NH3 (1,1) 23.6944955 GHz; NH3 (3,3) 23.8701296 GHz

\subsection*{Sources}
\texttt{JPL; SPLATALOGUE; SHIRLEY\_2015}

\clearpage


\section{SiO v=1 J=1-0 maser}
\index{SiO}

\begin{conceptbox}[At a glance]
\begin{tabularx}{\textwidth}{@{}lX@{}}
\textbf{Rest coordinate} & 43.122079 GHz \\
\textbf{Band} & millimetre / Q band \\
\textbf{Transition} & v=1, J=1 -> 0 \\
\textbf{Mechanism} & vibrationally excited rotational maser \\
\textbf{Primary tracers} & stellar-envelope dynamics; pulsation; magnetic-field morphology; compact shocks \\
\textbf{Typical environments} & inner envelopes of oxygen-rich AGB stars; red supergiants; rare shocked star-forming gas \\
\textbf{Indicative regime} & warm to hot molecular gas, often about 1000-2000 K near evolved stars \\
\textbf{Density guidance} & Requires high density and an inversion; representative pumping solutions use n(H2) around 10\textasciicircum{}9-10\textasciicircum{}10 cm\textasciicircum{}-3. \\
\end{tabularx}
\end{conceptbox}

SiO v=1 J=1-0 maser is produced by vibrationally excited rotational maser and is encountered in inner envelopes of oxygen-rich AGB stars; red supergiants; rare shocked star-forming gas. It is commonly used to trace stellar-envelope dynamics; pulsation; magnetic-field morphology; compact shocks.

\subsection*{How to use it}
Use this transition with its diagnostic companions to constrain stellar-envelope dynamics; pulsation; magnetic-field morphology; compact shocks, while fitting the continuum, line-spread function, and relevant blends.

\subsection*{Interpretive cautions}
Non-LTE, polarized and time variable. Do not use maser flux as a linear mass estimator.

\subsection*{Useful companions}
SiO v=2 J=1-0; SiO v=0 rotational lines; H2O 22 GHz maser

\subsection*{Sources}
\texttt{CDMS; JPL; SPLATALOGUE}

\clearpage


\section{HCN (1-0)}
\index{HCN}

\begin{conceptbox}[At a glance]
\begin{tabularx}{\textwidth}{@{}lX@{}}
\textbf{Rest coordinate} & 88.631847 GHz \\
\textbf{Band} & millimetre / 3 mm \\
\textbf{Transition} & J=1 -> 0; F=2 -> 1 strongest hyperfine component \\
\textbf{Mechanism} & electric-dipole rotational transition \\
\textbf{Primary tracers} & dense molecular gas; star-forming gas fraction; infall via self-absorption with isotopologues \\
\textbf{Typical environments} & dense molecular cores; starburst nuclei; AGN circumnuclear disks; molecular outflows \\
\textbf{Indicative regime} & commonly 10-100 K; IR pumping can matter in luminous dusty nuclei \\
\textbf{Density guidance} & Optically thin critical density is of order 10\textasciicircum{}6 cm\textasciicircum{}-3 near 10-20 K, while trapping makes effective excitation densities commonly 1-2 dex lower. \\
\end{tabularx}
\end{conceptbox}

HCN (1-0) is produced by electric-dipole rotational transition and is encountered in dense molecular cores; starburst nuclei; AGN circumnuclear disks; molecular outflows. It is commonly used to trace dense molecular gas; star-forming gas fraction; infall via self-absorption with isotopologues.

\subsection*{How to use it}
Use this transition with its diagnostic companions to constrain dense molecular gas; star-forming gas fraction; infall via self-absorption with isotopologues, while fitting the continuum, line-spread function, and relevant blends.

\subsection*{Interpretive cautions}
Resolve or model three hyperfine components in narrow Galactic sources. Opacity, abundance, electron excitation and 14-micron pumping complicate the dense-gas interpretation.

\subsection*{Useful companions}
HCO+ (1-0); HNC (1-0); H13CN (1-0); CN (1-0)

\subsection*{Sources}
\texttt{CDMS; SPLATALOGUE; SHIRLEY\_2015}

\clearpage


\section{HCO+ (1-0)}
\index{HCO+}

\begin{conceptbox}[At a glance]
\begin{tabularx}{\textwidth}{@{}lX@{}}
\textbf{Rest coordinate} & 89.188526 GHz \\
\textbf{Band} & millimetre / 3 mm \\
\textbf{Transition} & J=1 -> 0 \\
\textbf{Mechanism} & electric-dipole rotational transition of a molecular ion \\
\textbf{Primary tracers} & dense ionized-molecular chemistry; infall/outflow profiles; cosmic-ray ionization indirectly \\
\textbf{Typical environments} & dense molecular clouds; protostellar envelopes; galaxy nuclei; outflows \\
\textbf{Indicative regime} & about 10-100 K in most applications \\
\textbf{Density guidance} & Critical density is a few 10\textasciicircum{}5 cm\textasciicircum{}-3, but effective excitation density can be around 10\textasciicircum{}3-10\textasciicircum{}4 cm\textasciicircum{}-3 depending on opacity and column. \\
\end{tabularx}
\end{conceptbox}

HCO+ (1-0) is produced by electric-dipole rotational transition of a molecular ion and is encountered in dense molecular clouds; protostellar envelopes; galaxy nuclei; outflows. It is commonly used to trace dense ionized-molecular chemistry; infall/outflow profiles; cosmic-ray ionization indirectly.

\subsection*{How to use it}
Use this transition with its diagnostic companions to constrain dense ionized-molecular chemistry; infall/outflow profiles; cosmic-ray ionization indirectly, while fitting the continuum, line-spread function, and relevant blends.

\subsection*{Interpretive cautions}
Often optically thick and self-absorbed. Pair with H13CO+ to estimate opacity; abundance is sensitive to ionization and water/CO chemistry.

\subsection*{Useful companions}
HCN (1-0); H13CO+ (1-0); CO isotopologues; N2H+ (1-0)

\subsection*{Sources}
\texttt{CDMS; SPLATALOGUE; SHIRLEY\_2015}

\clearpage


\section{N2H+ (1-0)}
\index{N2H+}

\begin{conceptbox}[At a glance]
\begin{tabularx}{\textwidth}{@{}lX@{}}
\textbf{Rest coordinate} & 93.173764 GHz \\
\textbf{Band} & millimetre / 3 mm \\
\textbf{Transition} & J=1 -> 0 hyperfine multiplet \\
\textbf{Mechanism} & electric-dipole rotational transition with nitrogen hyperfine structure \\
\textbf{Primary tracers} & cold dense gas; systemic velocity; optical depth from hyperfine fitting; CO freeze-out zones \\
\textbf{Typical environments} & cold prestellar cores; dense filaments; protostellar envelopes; CO-depleted gas \\
\textbf{Indicative regime} & especially useful around 7-30 K \\
\textbf{Density guidance} & Critical density is roughly 10\textasciicircum{}5 cm\textasciicircum{}-3; radiative trapping and hyperfine excitation lower the effective threshold. \\
\end{tabularx}
\end{conceptbox}

N2H+ (1-0) is produced by electric-dipole rotational transition with nitrogen hyperfine structure and is encountered in cold prestellar cores; dense filaments; protostellar envelopes; CO-depleted gas. It is commonly used to trace cold dense gas; systemic velocity; optical depth from hyperfine fitting; CO freeze-out zones.

\subsection*{How to use it}
Use this transition with its diagnostic companions to constrain cold dense gas; systemic velocity; optical depth from hyperfine fitting; CO freeze-out zones, while fitting the continuum, line-spread function, and relevant blends.

\subsection*{Interpretive cautions}
The quoted value is a conventional reference for the hyperfine complex; precision velocity work must adopt component-specific frequencies and a consistent catalog.

\subsection*{Useful companions}
NH3 (1,1); HCO+ (1-0); C18O (1-0); N2D+

\subsection*{Sources}
\texttt{CDMS; SPLATALOGUE; SHIRLEY\_2015}

\clearpage


\section{CS (2-1)}
\index{CS}

\begin{conceptbox}[At a glance]
\begin{tabularx}{\textwidth}{@{}lX@{}}
\textbf{Rest coordinate} & 97.980953 GHz \\
\textbf{Band} & millimetre / 3 mm \\
\textbf{Transition} & J=2 -> 1 \\
\textbf{Mechanism} & electric-dipole rotational transition \\
\textbf{Primary tracers} & dense gas; velocity structure; sulfur chemistry \\
\textbf{Typical environments} & dense clumps; massive star-forming regions; galactic centers; shocked molecular gas \\
\textbf{Indicative regime} & about 10-100 K \\
\textbf{Density guidance} & Critical density is a few 10\textasciicircum{}5 cm\textasciicircum{}-3, with a substantially lower effective excitation density for optically thick bright columns. \\
\end{tabularx}
\end{conceptbox}

CS (2-1) is produced by electric-dipole rotational transition and is encountered in dense clumps; massive star-forming regions; galactic centers; shocked molecular gas. It is commonly used to trace dense gas; velocity structure; sulfur chemistry.

\subsection*{How to use it}
Use this transition with its diagnostic companions to constrain dense gas; velocity structure; sulfur chemistry, while fitting the continuum, line-spread function, and relevant blends.

\subsection*{Interpretive cautions}
CS can deplete in cold cores and become optically thick; use isotopologues and multi-J excitation rather than a single conversion factor.

\subsection*{Useful companions}
C34S (2-1); HCN (1-0); CS higher-J ladder

\subsection*{Sources}
\texttt{CDMS; JPL; SHIRLEY\_2015}

\clearpage


\section{C18O (1-0)}
\index{C18O}

\begin{conceptbox}[At a glance]
\begin{tabularx}{\textwidth}{@{}lX@{}}
\textbf{Rest coordinate} & 109.7821734 GHz \\
\textbf{Band} & millimetre / 3 mm \\
\textbf{Transition} & J=1 -> 0 \\
\textbf{Mechanism} & electric-dipole rotational transition \\
\textbf{Primary tracers} & moderate-column molecular gas; CO optical-depth correction; gas kinematics \\
\textbf{Typical environments} & molecular clouds; dense cores; protoplanetary disks; galaxy molecular disks \\
\textbf{Indicative regime} & cold molecular gas, commonly 8-50 K \\
\textbf{Density guidance} & Critical density is around 10\textasciicircum{}3 cm\textasciicircum{}-3 and effective excitation density lower; abundance and shielding often dominate detectability. \\
\end{tabularx}
\end{conceptbox}

C18O (1-0) is produced by electric-dipole rotational transition and is encountered in molecular clouds; dense cores; protoplanetary disks; galaxy molecular disks. It is commonly used to trace moderate-column molecular gas; CO optical-depth correction; gas kinematics.

\subsection*{How to use it}
Use this transition with its diagnostic companions to constrain moderate-column molecular gas; CO optical-depth correction; gas kinematics, while fitting the continuum, line-spread function, and relevant blends.

\subsection*{Interpretive cautions}
Usually thinner than 13CO but not guaranteed optically thin. Selective photodissociation, fractionation and isotope-ratio gradients bias mass estimates.

\subsection*{Useful companions}
13CO (1-0); 12CO (1-0); C17O (1-0); dust continuum

\subsection*{Sources}
\texttt{CDMS; JPL; CARILLI\_WALTER\_2013}

\clearpage


\section{13CO (1-0)}
\index{13CO}

\begin{conceptbox}[At a glance]
\begin{tabularx}{\textwidth}{@{}lX@{}}
\textbf{Rest coordinate} & 110.2013543 GHz \\
\textbf{Band} & millimetre / 3 mm \\
\textbf{Transition} & J=1 -> 0 \\
\textbf{Mechanism} & electric-dipole rotational transition \\
\textbf{Primary tracers} & molecular column; 12CO opacity; cloud structure; isotope chemistry \\
\textbf{Typical environments} & molecular clouds; galaxy disks; circumstellar envelopes; protoplanetary disks \\
\textbf{Indicative regime} & commonly 8-100 K \\
\textbf{Density guidance} & Critical density is around 10\textasciicircum{}3 cm\textasciicircum{}-3; photon trapping lowers the effective excitation threshold. \\
\end{tabularx}
\end{conceptbox}

13CO (1-0) is produced by electric-dipole rotational transition and is encountered in molecular clouds; galaxy disks; circumstellar envelopes; protoplanetary disks. It is commonly used to trace molecular column; 12CO opacity; cloud structure; isotope chemistry.

\subsection*{How to use it}
Use this transition with its diagnostic companions to constrain molecular column; 12CO opacity; cloud structure; isotope chemistry, while fitting the continuum, line-spread function, and relevant blends.

\subsection*{Interpretive cautions}
Often optically thick in cloud peaks. Carbon isotope gradients, chemical fractionation and selective photodissociation prevent a universal 12CO/13CO abundance ratio.

\subsection*{Useful companions}
12CO (1-0); C18O (1-0); 13CO higher-J lines

\subsection*{Sources}
\texttt{CDMS; JPL; CARILLI\_WALTER\_2013}

\clearpage


\section{12CO (1-0)}
\index{12CO}

\begin{conceptbox}[At a glance]
\begin{tabularx}{\textwidth}{@{}lX@{}}
\textbf{Rest coordinate} & 115.2712018 GHz \\
\textbf{Band} & millimetre / 2.6 mm \\
\textbf{Transition} & J=1 -> 0 \\
\textbf{Mechanism} & electric-dipole rotational transition \\
\textbf{Primary tracers} & bulk molecular gas via CO-to-H2 conversion; kinematics; cloud temperature when thick \\
\textbf{Typical environments} & molecular clouds; galaxy disks; starbursts; circumstellar envelopes; molecular outflows \\
\textbf{Indicative regime} & cold molecular gas from about 5 K upward; upper-state energy is 5.53 K \\
\textbf{Density guidance} & Optically thin critical density is roughly 2x10\textasciicircum{}3 cm\textasciicircum{}-3, while trapping permits bright emission at a few 10\textasciicircum{}2 cm\textasciicircum{}-3 or below. \\
\end{tabularx}
\end{conceptbox}

12CO (1-0) is produced by electric-dipole rotational transition and is encountered in molecular clouds; galaxy disks; starbursts; circumstellar envelopes; molecular outflows. It is commonly used to trace bulk molecular gas via CO-to-H2 conversion; kinematics; cloud temperature when thick.

\subsection*{How to use it}
Use this transition with its diagnostic companions to constrain bulk molecular gas via CO-to-H2 conversion; kinematics; cloud temperature when thick, while fitting the continuum, line-spread function, and relevant blends.

\subsection*{Interpretive cautions}
Usually optically thick. X\_CO depends on metallicity, dynamics, excitation, spatial scale and radiation field; CMB contrast matters at high redshift.

\subsection*{Useful companions}
13CO (1-0); C18O (1-0); CO (2-1); dust continuum

\subsection*{Sources}
\texttt{CDMS; JPL; CARILLI\_WALTER\_2013}

\clearpage


\section{12CO (2-1)}
\index{12CO}

\begin{conceptbox}[At a glance]
\begin{tabularx}{\textwidth}{@{}lX@{}}
\textbf{Rest coordinate} & 230.538 GHz \\
\textbf{Band} & submillimetre / 1.3 mm \\
\textbf{Transition} & J=2 -> 1 \\
\textbf{Mechanism} & electric-dipole rotational transition \\
\textbf{Primary tracers} & molecular gas distribution; excitation with CO(1-0); kinematics \\
\textbf{Typical environments} & molecular clouds; nearby and high-redshift galaxies; disks; outflows \\
\textbf{Indicative regime} & cold to moderately warm gas; upper-state energy is 16.6 K \\
\textbf{Density guidance} & Critical density is around 10\textasciicircum{}4 cm\textasciicircum{}-3, but opacity makes the effective density much lower. \\
\end{tabularx}
\end{conceptbox}

12CO (2-1) is produced by electric-dipole rotational transition and is encountered in molecular clouds; nearby and high-redshift galaxies; disks; outflows. It is commonly used to trace molecular gas distribution; excitation with CO(1-0); kinematics.

\subsection*{How to use it}
Use this transition with its diagnostic companions to constrain molecular gas distribution; excitation with CO(1-0); kinematics, while fitting the continuum, line-spread function, and relevant blends.

\subsection*{Interpretive cautions}
Brightness-temperature ratio r21 is sensitive to excitation, beam filling and optical depth. Match angular resolution before interpreting line ratios.

\subsection*{Useful companions}
CO (1-0); CO (3-2); 13CO (2-1)

\subsection*{Sources}
\texttt{CDMS; JPL; CARILLI\_WALTER\_2013}

\clearpage


\section{12CO (3-2)}
\index{12CO}

\begin{conceptbox}[At a glance]
\begin{tabularx}{\textwidth}{@{}lX@{}}
\textbf{Rest coordinate} & 345.7959899 GHz \\
\textbf{Band} & submillimetre / 870 micron \\
\textbf{Transition} & J=3 -> 2 \\
\textbf{Mechanism} & electric-dipole rotational transition \\
\textbf{Primary tracers} & warmer/denser molecular component; CO excitation ladder; kinematics \\
\textbf{Typical environments} & warm molecular clouds; star-forming clumps; starbursts; high-redshift galaxies \\
\textbf{Indicative regime} & commonly 20-100 K; upper-state energy is 33.2 K \\
\textbf{Density guidance} & Critical density is a few 10\textasciicircum{}4 cm\textasciicircum{}-3; radiative trapping lowers the effective density. \\
\end{tabularx}
\end{conceptbox}

12CO (3-2) is produced by electric-dipole rotational transition and is encountered in warm molecular clouds; star-forming clumps; starbursts; high-redshift galaxies. It is commonly used to trace warmer/denser molecular component; CO excitation ladder; kinematics.

\subsection*{How to use it}
Use this transition with its diagnostic companions to constrain warmer/denser molecular component; CO excitation ladder; kinematics, while fitting the continuum, line-spread function, and relevant blends.

\subsection*{Interpretive cautions}
Mid-J selection can bias samples toward excited gas. Use a CO SLED and consistent beam/uv coverage to infer multiple gas phases.

\subsection*{Useful companions}
CO (2-1); CO (4-3); 13CO (3-2); [C I] 492 GHz

\subsection*{Sources}
\texttt{CDMS; JPL; CARILLI\_WALTER\_2013}

\clearpage


\section{[C I] 609 micron}
\index{C I}

\begin{conceptbox}[At a glance]
\begin{tabularx}{\textwidth}{@{}lX@{}}
\textbf{Rest coordinate} & 492.160651 GHz \\
\textbf{Band} & submillimetre / far-infrared \\
\textbf{Transition} & 3P1 -> 3P0 \\
\textbf{Mechanism} & magnetic-dipole fine-structure transition \\
\textbf{Primary tracers} & neutral carbon mass; molecular gas when calibrated; excitation with [C I](2-1) \\
\textbf{Typical environments} & molecular clouds; PDRs; high-redshift galaxies; CO-dark molecular interfaces \\
\textbf{Indicative regime} & cold molecular gas, commonly 15-60 K; upper level is 23.6 K above ground \\
\textbf{Density guidance} & Critical density with H2 is of order 5x10\textasciicircum{}2 cm\textasciicircum{}-3, weakly model dependent. \\
\end{tabularx}
\end{conceptbox}

[C I] 609 micron is produced by magnetic-dipole fine-structure transition and is encountered in molecular clouds; PDRs; high-redshift galaxies; CO-dark molecular interfaces. It is commonly used to trace neutral carbon mass; molecular gas when calibrated; excitation with [C I](2-1).

\subsection*{How to use it}
Use this transition with its diagnostic companions to constrain neutral carbon mass; molecular gas when calibrated; excitation with [C I](2-1), while fitting the continuum, line-spread function, and relevant blends.

\subsection*{Interpretive cautions}
The two [C I] lines give an excitation temperature only under a common, often optically thin one-zone assumption. Carbon abundance is not universal.

\subsection*{Useful companions}
[C I] 809 GHz; CO rotational ladder; [C II] 158 micron

\subsection*{Sources}
\texttt{NIST\_ASD; CDMS; CARILLI\_WALTER\_2013}

\clearpage


\section{[C I] 370 micron}
\index{C I}

\begin{conceptbox}[At a glance]
\begin{tabularx}{\textwidth}{@{}lX@{}}
\textbf{Rest coordinate} & 809.34197 GHz \\
\textbf{Band} & submillimetre / far-infrared \\
\textbf{Transition} & 3P2 -> 3P1 \\
\textbf{Mechanism} & magnetic-dipole fine-structure transition \\
\textbf{Primary tracers} & neutral carbon excitation; molecular gas temperature; carbon mass \\
\textbf{Typical environments} & warm molecular clouds; PDRs; starbursts; high-redshift galaxies \\
\textbf{Indicative regime} & typically 25-100 K; upper level is 62.5 K above ground \\
\textbf{Density guidance} & Critical density with H2 is of order 10\textasciicircum{}3 cm\textasciicircum{}-3. \\
\end{tabularx}
\end{conceptbox}

[C I] 370 micron is produced by magnetic-dipole fine-structure transition and is encountered in warm molecular clouds; PDRs; starbursts; high-redshift galaxies. It is commonly used to trace neutral carbon excitation; molecular gas temperature; carbon mass.

\subsection*{How to use it}
Use this transition with its diagnostic companions to constrain neutral carbon excitation; molecular gas temperature; carbon mass, while fitting the continuum, line-spread function, and relevant blends.

\subsection*{Interpretive cautions}
The line sits close to CO(7-6) in frequency, useful for simultaneous setups but demanding correct identification at uncertain redshift.

\subsection*{Useful companions}
[C I] 492 GHz; CO (7-6) nearby in frequency; [C II] 158 micron

\subsection*{Sources}
\texttt{NIST\_ASD; CDMS; CARILLI\_WALTER\_2013}

\clearpage


\section{[C II] 158 micron}
\index{C II}

\begin{conceptbox}[At a glance]
\begin{tabularx}{\textwidth}{@{}lX@{}}
\textbf{Rest coordinate} & 1900.5369 GHz \\
\textbf{Band} & far-infrared \\
\textbf{Transition} & 2P3/2 -> 2P1/2 \\
\textbf{Mechanism} & fine-structure transition \\
\textbf{Primary tracers} & gas cooling; star-forming ISM; dynamics and redshift; CO-dark molecular gas with modeling \\
\textbf{Typical environments} & PDR surfaces; diffuse neutral gas; CO-dark H2; diffuse ionized gas; high-redshift galaxies \\
\textbf{Indicative regime} & upper level corresponds to 91.2 K; excited in neutral and ionized gas from roughly 50 K to 10\textasciicircum{}4 K \\
\textbf{Density guidance} & Approximate critical densities: about 50 cm\textasciicircum{}-3 for electrons, 3x10\textasciicircum{}3 for H and 5x10\textasciicircum{}3 for H2, with temperature dependence. \\
\end{tabularx}
\end{conceptbox}

[C II] 158 micron is produced by fine-structure transition and is encountered in PDR surfaces; diffuse neutral gas; CO-dark H2; diffuse ionized gas; high-redshift galaxies. It is commonly used to trace gas cooling; star-forming ISM; dynamics and redshift; CO-dark molecular gas with modeling.

\subsection*{How to use it}
Use this transition with its diagnostic companions to constrain gas cooling; star-forming ISM; dynamics and redshift; CO-dark molecular gas with modeling, while fitting the continuum, line-spread function, and relevant blends.

\subsection*{Interpretive cautions}
A multiphase line, not a pure PDR or SFR tracer. Optical depth, foreground absorption and the high-surface-brightness [C II] deficit affect interpretation.

\subsection*{Useful companions}
[N II] 205 micron to isolate ionized contribution; [O I] 63 micron; CO rotational lines; FIR continuum

\subsection*{Sources}
\texttt{NIST\_ASD; CARILLI\_WALTER\_2013; HOLLENBACH\_TIELENS\_1999}

\clearpage


\section{[N II] 205 micron}
\index{N II}

\begin{conceptbox}[At a glance]
\begin{tabularx}{\textwidth}{@{}lX@{}}
\textbf{Rest coordinate} & 1461.1319 GHz \\
\textbf{Band} & far-infrared \\
\textbf{Transition} & 3P1 -> 3P0 \\
\textbf{Mechanism} & fine-structure transition \\
\textbf{Primary tracers} & ionized-gas mass; massive-star ionization; ionized contribution to [C II] \\
\textbf{Typical environments} & diffuse ionized medium; H II regions; star-forming galaxies \\
\textbf{Indicative regime} & ionized gas near 6000-10000 K; excitation energy is only about 70 K \\
\textbf{Density guidance} & Electron critical density is roughly 40-50 cm\textasciicircum{}-3; the 122/205 ratio measures electron density over roughly 1-10\textasciicircum{}3 cm\textasciicircum{}-3. \\
\end{tabularx}
\end{conceptbox}

[N II] 205 micron is produced by fine-structure transition and is encountered in diffuse ionized medium; H II regions; star-forming galaxies. It is commonly used to trace ionized-gas mass; massive-star ionization; ionized contribution to [C II].

\subsection*{How to use it}
Use this transition with its diagnostic companions to constrain ionized-gas mass; massive-star ionization; ionized contribution to [C II], while fitting the continuum, line-spread function, and relevant blends.

\subsection*{Interpretive cautions}
N+ requires ionized gas because its formation threshold exceeds 13.6 eV. Low critical density makes 205 micron particularly sensitive to diffuse ionized material.

\subsection*{Useful companions}
[N II] 122 micron; [C II] 158 micron; radio free-free emission

\subsection*{Sources}
\texttt{NIST\_NII; CARILLI\_WALTER\_2013; OSTERBROCK\_FERLAND\_2006}

\clearpage


\section{[N II] 122 micron}
\index{N II}

\begin{conceptbox}[At a glance]
\begin{tabularx}{\textwidth}{@{}lX@{}}
\textbf{Rest coordinate} & 2459.38006 GHz \\
\textbf{Band} & far-infrared \\
\textbf{Transition} & 3P2 -> 3P1 \\
\textbf{Mechanism} & fine-structure transition \\
\textbf{Primary tracers} & electron density with 205 micron; ionized-gas cooling; massive-star photoionization \\
\textbf{Typical environments} & H II regions; diffuse ionized gas; star-forming galaxies \\
\textbf{Indicative regime} & ionized gas around 6000-10000 K \\
\textbf{Density guidance} & Electron critical density is about 300 cm\textasciicircum{}-3; 122/205 is a robust density ratio when both lines sample the same gas. \\
\end{tabularx}
\end{conceptbox}

[N II] 122 micron is produced by fine-structure transition and is encountered in H II regions; diffuse ionized gas; star-forming galaxies. It is commonly used to trace electron density with 205 micron; ionized-gas cooling; massive-star photoionization.

\subsection*{How to use it}
Use this transition with its diagnostic companions to constrain electron density with 205 micron; ionized-gas cooling; massive-star photoionization, while fitting the continuum, line-spread function, and relevant blends.

\subsection*{Interpretive cautions}
Far-IR dust extinction is small, but aperture mismatch and multi-density structure can bias the ratio.

\subsection*{Useful companions}
[N II] 205 micron; [O III] 88 micron; [C II] 158 micron

\subsection*{Sources}
\texttt{NIST\_NII; CARILLI\_WALTER\_2013; OSTERBROCK\_FERLAND\_2006}

\clearpage


\section{[O I] 145 micron}
\index{O I}

\begin{conceptbox}[At a glance]
\begin{tabularx}{\textwidth}{@{}lX@{}}
\textbf{Rest coordinate} & 145.525 micron \\
\textbf{Band} & far-infrared \\
\textbf{Transition} & 3P0 -> 3P1 \\
\textbf{Mechanism} & fine-structure transition \\
\textbf{Primary tracers} & neutral-gas cooling; PDR density and radiation field with [O I] 63/[C II]; shock-heated gas \\
\textbf{Typical environments} & warm dense PDRs; shocks; neutral gas near H II regions \\
\textbf{Indicative regime} & warm neutral gas, commonly about 100-1000 K \\
\textbf{Density guidance} & Critical density for H/H2 collisions is of order 10\textasciicircum{}5 cm\textasciicircum{}-3, model dependent. \\
\end{tabularx}
\end{conceptbox}

[O I] 145 micron is produced by fine-structure transition and is encountered in warm dense PDRs; shocks; neutral gas near H II regions. It is commonly used to trace neutral-gas cooling; PDR density and radiation field with [O I] 63/[C II]; shock-heated gas.

\subsection*{How to use it}
Use this transition with its diagnostic companions to constrain neutral-gas cooling; PDR density and radiation field with [O I] 63/[C II]; shock-heated gas, while fitting the continuum, line-spread function, and relevant blends.

\subsection*{Interpretive cautions}
Usually weaker but less self-absorbed than the 63-micron line. The 63/145 ratio can reveal opacity, foreground absorption, temperature and density.

\subsection*{Useful companions}
[O I] 63 micron; [C II] 158 micron; high-J CO

\subsection*{Sources}
\texttt{NIST\_ASD; HOLLENBACH\_TIELENS\_1999; CARILLI\_WALTER\_2013}

\clearpage


\section{[O I] 63 micron}
\index{O I}

\begin{conceptbox}[At a glance]
\begin{tabularx}{\textwidth}{@{}lX@{}}
\textbf{Rest coordinate} & 63.1837 micron \\
\textbf{Band} & far-infrared \\
\textbf{Transition} & 3P1 -> 3P2 \\
\textbf{Mechanism} & fine-structure transition \\
\textbf{Primary tracers} & dominant neutral-gas cooling in warm dense regions; FUV field; shock power \\
\textbf{Typical environments} & dense warm PDRs; dissociative and non-dissociative shocks; protostellar outflows \\
\textbf{Indicative regime} & upper level is about 228 K; strongest in gas roughly 100-1000 K \\
\textbf{Density guidance} & Critical density for collisions with H is roughly 5x10\textasciicircum{}5 cm\textasciicircum{}-3; optical depth often matters first. \\
\end{tabularx}
\end{conceptbox}

[O I] 63 micron is produced by fine-structure transition and is encountered in dense warm PDRs; dissociative and non-dissociative shocks; protostellar outflows. It is commonly used to trace dominant neutral-gas cooling in warm dense regions; FUV field; shock power.

\subsection*{How to use it}
Use this transition with its diagnostic companions to constrain dominant neutral-gas cooling in warm dense regions; FUV field; shock power, while fitting the continuum, line-spread function, and relevant blends.

\subsection*{Interpretive cautions}
Can be optically thick and self-absorbed by colder foreground O I; a line deficit does not uniquely imply weak PDR heating.

\subsection*{Useful companions}
[O I] 145 micron; [C II] 158 micron; H2 rotational lines

\subsection*{Sources}
\texttt{NIST\_ASD; HOLLENBACH\_TIELENS\_1999; CARILLI\_WALTER\_2013}

\clearpage


\section{[O III] 88 micron}
\index{O III}

\begin{conceptbox}[At a glance]
\begin{tabularx}{\textwidth}{@{}lX@{}}
\textbf{Rest coordinate} & 3393.006 GHz \\
\textbf{Band} & far-infrared \\
\textbf{Transition} & 3P1 -> 3P0 \\
\textbf{Mechanism} & fine-structure transition \\
\textbf{Primary tracers} & hard massive-star radiation; ionized gas; electron density with 52 micron; high-redshift systemic redshift \\
\textbf{Typical environments} & H II regions; low-metallicity star-forming galaxies; planetary nebulae; reionization-era galaxies \\
\textbf{Indicative regime} & photoionized gas around 7000-20000 K; low excitation energy minimizes temperature sensitivity \\
\textbf{Density guidance} & Electron critical density is about 500 cm\textasciicircum{}-3; 52/88 measures densities from roughly 10\textasciicircum{}2 to 10\textasciicircum{}4 cm\textasciicircum{}-3. \\
\end{tabularx}
\end{conceptbox}

[O III] 88 micron is produced by fine-structure transition and is encountered in H II regions; low-metallicity star-forming galaxies; planetary nebulae; reionization-era galaxies. It is commonly used to trace hard massive-star radiation; ionized gas; electron density with 52 micron; high-redshift systemic redshift.

\subsection*{How to use it}
Use this transition with its diagnostic companions to constrain hard massive-star radiation; ionized gas; electron density with 52 micron; high-redshift systemic redshift, while fitting the continuum, line-spread function, and relevant blends.

\subsection*{Interpretive cautions}
Bright in hard, low-metallicity H II regions and relatively extinction-free. Abundance, ionization parameter and density all affect line-to-SFR conversions.

\subsection*{Useful companions}
[O III] 52 micron; [N II] 122 micron; optical [O III] 5008 A

\subsection*{Sources}
\texttt{NIST\_ASD; CARILLI\_WALTER\_2013; OSTERBROCK\_FERLAND\_2006}

\clearpage


\section{[O III] 52 micron}
\index{O III}

\begin{conceptbox}[At a glance]
\begin{tabularx}{\textwidth}{@{}lX@{}}
\textbf{Rest coordinate} & 51.8145 micron \\
\textbf{Band} & far-infrared \\
\textbf{Transition} & 3P2 -> 3P1 \\
\textbf{Mechanism} & fine-structure transition \\
\textbf{Primary tracers} & electron density with 88 micron; hard ionizing field; ionized-gas cooling \\
\textbf{Typical environments} & compact H II regions; planetary nebulae; starbursts \\
\textbf{Indicative regime} & photoionized gas around 7000-20000 K \\
\textbf{Density guidance} & Electron critical density is roughly 3.6x10\textasciicircum{}3 cm\textasciicircum{}-3; 52/88 is a standard far-IR density diagnostic. \\
\end{tabularx}
\end{conceptbox}

[O III] 52 micron is produced by fine-structure transition and is encountered in compact H II regions; planetary nebulae; starbursts. It is commonly used to trace electron density with 88 micron; hard ionizing field; ionized-gas cooling.

\subsection*{How to use it}
Use this transition with its diagnostic companions to constrain electron density with 88 micron; hard ionizing field; ionized-gas cooling, while fitting the continuum, line-spread function, and relevant blends.

\subsection*{Interpretive cautions}
Use matched apertures. At high density both lines approach LTE and the ratio saturates; abundance and ionization structure remain separate uncertainties.

\subsection*{Useful companions}
[O III] 88 micron; optical [O III] 5008 A; [N III] 57 micron

\subsection*{Sources}
\texttt{MPE\_FINE\_STRUCTURE; CARILLI\_WALTER\_2013; OSTERBROCK\_FERLAND\_2006}

\clearpage


\section{[Si II] 34.8 micron}
\index{Si II}

\begin{conceptbox}[At a glance]
\begin{tabularx}{\textwidth}{@{}lX@{}}
\textbf{Rest coordinate} & 34.814 micron \\
\textbf{Band} & mid-infrared \\
\textbf{Transition} & 2P3/2 -> 2P1/2 \\
\textbf{Mechanism} & fine-structure transition \\
\textbf{Primary tracers} & gas cooling; silicon released from grains; PDR/XDR/shock conditions \\
\textbf{Typical environments} & PDRs; diffuse ionized gas; shocks; X-ray-dominated regions \\
\textbf{Indicative regime} & warm neutral or ionized gas, about 100-10000 K \\
\textbf{Density guidance} & Critical density is roughly 10\textasciicircum{}3 cm\textasciicircum{}-3 for electrons and a few 10\textasciicircum{}5 cm\textasciicircum{}-3 for H, with temperature dependence. \\
\end{tabularx}
\end{conceptbox}

[Si II] 34.8 micron is produced by fine-structure transition and is encountered in PDRs; diffuse ionized gas; shocks; X-ray-dominated regions. It is commonly used to trace gas cooling; silicon released from grains; PDR/XDR/shock conditions.

\subsection*{How to use it}
Use this transition with its diagnostic companions to constrain gas cooling; silicon released from grains; PDR/XDR/shock conditions, while fitting the continuum, line-spread function, and relevant blends.

\subsection*{Interpretive cautions}
Silicon depletion and grain destruction dominate absolute calibration; strong emission may arise from multiple phases along the same sightline.

\subsection*{Useful companions}
[Fe II] 26 micron; [S III] 33.5 micron; [C II] 158 micron

\subsection*{Sources}
\texttt{ISO\_SWS; MPE\_FINE\_STRUCTURE; HOLLENBACH\_TIELENS\_1999}

\clearpage


\section{[S III] 33.5 micron}
\index{S III}

\begin{conceptbox}[At a glance]
\begin{tabularx}{\textwidth}{@{}lX@{}}
\textbf{Rest coordinate} & 33.481 micron \\
\textbf{Band} & mid-infrared \\
\textbf{Transition} & 3P1 -> 3P0 \\
\textbf{Mechanism} & fine-structure transition \\
\textbf{Primary tracers} & ionized gas; electron density with 18.7 micron; sulfur abundance with modeling \\
\textbf{Typical environments} & H II regions; starburst galaxies; planetary nebulae \\
\textbf{Indicative regime} & photoionized gas around 5000-15000 K \\
\textbf{Density guidance} & Electron critical density is of order 10\textasciicircum{}3 cm\textasciicircum{}-3; the 18.7/33.5 ratio spans roughly 10\textasciicircum{}2-10\textasciicircum{}4 cm\textasciicircum{}-3. \\
\end{tabularx}
\end{conceptbox}

[S III] 33.5 micron is produced by fine-structure transition and is encountered in H II regions; starburst galaxies; planetary nebulae. It is commonly used to trace ionized gas; electron density with 18.7 micron; sulfur abundance with modeling.

\subsection*{How to use it}
Use this transition with its diagnostic companions to constrain ionized gas; electron density with 18.7 micron; sulfur abundance with modeling, while fitting the continuum, line-spread function, and relevant blends.

\subsection*{Interpretive cautions}
The widely separated pair is susceptible to aperture and extinction mismatch; use the correct low-density limit for the adopted atomic data.

\subsection*{Useful companions}
[S III] 18.7 micron; optical [S III] 9071/9533 A; [Ne II] 12.8 micron

\subsection*{Sources}
\texttt{ISO\_SWS; MPE\_FINE\_STRUCTURE; OSTERBROCK\_FERLAND\_2006}

\clearpage


\section{[S III] 18.7 micron}
\index{S III}

\begin{conceptbox}[At a glance]
\begin{tabularx}{\textwidth}{@{}lX@{}}
\textbf{Rest coordinate} & 18.713 micron \\
\textbf{Band} & mid-infrared \\
\textbf{Transition} & 3P2 -> 3P1 \\
\textbf{Mechanism} & fine-structure transition \\
\textbf{Primary tracers} & electron density with 33.5 micron; ionized gas; obscured star formation \\
\textbf{Typical environments} & H II regions; compact starbursts; planetary nebulae \\
\textbf{Indicative regime} & photoionized gas around 5000-15000 K \\
\textbf{Density guidance} & Electron critical density is roughly 2x10\textasciicircum{}4 cm\textasciicircum{}-3; 18.7/33.5 rises with density before saturating. \\
\end{tabularx}
\end{conceptbox}

[S III] 18.7 micron is produced by fine-structure transition and is encountered in H II regions; compact starbursts; planetary nebulae. It is commonly used to trace electron density with 33.5 micron; ionized gas; obscured star formation.

\subsection*{How to use it}
Use this transition with its diagnostic companions to constrain electron density with 33.5 micron; ionized gas; obscured star formation, while fitting the continuum, line-spread function, and relevant blends.

\subsection*{Interpretive cautions}
Less affected by dust than optical lines but can suffer silicate-feature continuum systematics and aperture differences.

\subsection*{Useful companions}
[S III] 33.5 micron; [Ne III] 15.6 micron; optical [S II] density pair

\subsection*{Sources}
\texttt{ISO\_SWS; MPE\_FINE\_STRUCTURE; OSTERBROCK\_FERLAND\_2006}

\clearpage


\section{[Ne II] 12.8 micron}
\index{Ne II}

\begin{conceptbox}[At a glance]
\begin{tabularx}{\textwidth}{@{}lX@{}}
\textbf{Rest coordinate} & 12.814 micron \\
\textbf{Band} & mid-infrared \\
\textbf{Transition} & 2P1/2 -> 2P3/2 \\
\textbf{Mechanism} & fine-structure transition \\
\textbf{Primary tracers} & obscured massive-star ionization; ionized-gas kinematics; soft-to-moderate radiation fields \\
\textbf{Typical environments} & H II regions; embedded star formation; protoplanetary disk surfaces; AGN narrow-line regions \\
\textbf{Indicative regime} & photoionized gas around 5000-15000 K \\
\textbf{Density guidance} & Electron critical density is of order 6x10\textasciicircum{}5 cm\textasciicircum{}-3; usually in the low-density limit for extended H II gas. \\
\end{tabularx}
\end{conceptbox}

[Ne II] 12.8 micron is produced by fine-structure transition and is encountered in H II regions; embedded star formation; protoplanetary disk surfaces; AGN narrow-line regions. It is commonly used to trace obscured massive-star ionization; ionized-gas kinematics; soft-to-moderate radiation fields.

\subsection*{How to use it}
Use this transition with its diagnostic companions to constrain obscured massive-star ionization; ionized-gas kinematics; soft-to-moderate radiation fields, while fitting the continuum, line-spread function, and relevant blends.

\subsection*{Interpretive cautions}
[Ne III]/[Ne II] is sensitive to ionizing hardness and ionization parameter, but also geometry, metallicity and stellar population age.

\subsection*{Useful companions}
[Ne III] 15.6 micron; hydrogen recombination lines; [S IV] 10.5 micron

\subsection*{Sources}
\texttt{ISO\_SWS; MPE\_FINE\_STRUCTURE; KENNICUTT\_EVANS\_2012}

\clearpage


\section{[Ne III] 15.6 micron}
\index{Ne III}

\begin{conceptbox}[At a glance]
\begin{tabularx}{\textwidth}{@{}lX@{}}
\textbf{Rest coordinate} & 15.5551 micron \\
\textbf{Band} & mid-infrared \\
\textbf{Transition} & 3P1 -> 3P2 \\
\textbf{Mechanism} & fine-structure transition \\
\textbf{Primary tracers} & hard ionizing radiation; obscured star formation; ionization parameter with other neon lines \\
\textbf{Typical environments} & high-excitation H II regions; low-metallicity starbursts; planetary nebulae; AGN narrow-line regions \\
\textbf{Indicative regime} & photoionized gas around 7000-20000 K \\
\textbf{Density guidance} & Electron critical density is about 10\textasciicircum{}5 cm\textasciicircum{}-3; typically low-density in integrated galaxies. \\
\end{tabularx}
\end{conceptbox}

[Ne III] 15.6 micron is produced by fine-structure transition and is encountered in high-excitation H II regions; low-metallicity starbursts; planetary nebulae; AGN narrow-line regions. It is commonly used to trace hard ionizing radiation; obscured star formation; ionization parameter with other neon lines.

\subsection*{How to use it}
Use this transition with its diagnostic companions to constrain hard ionizing radiation; obscured star formation; ionization parameter with other neon lines, while fitting the continuum, line-spread function, and relevant blends.

\subsection*{Interpretive cautions}
Neon is only weakly depleted, but the ratio to [Ne II] is not a one-parameter hardness indicator.

\subsection*{Useful companions}
[Ne II] 12.8 micron; [Ne V] 14.3 micron; optical [Ne III] 3869.9 A

\subsection*{Sources}
\texttt{ISO\_SWS; MPE\_FINE\_STRUCTURE; KENNICUTT\_EVANS\_2012}

\clearpage


\section{H2 0-0 S(0)}
\index{H2}

\begin{conceptbox}[At a glance]
\begin{tabularx}{\textwidth}{@{}lX@{}}
\textbf{Rest coordinate} & 28.21883 micron \\
\textbf{Band} & mid-infrared \\
\textbf{Transition} & v=0, J=2 -> 0 (para-H2) \\
\textbf{Mechanism} & electric-quadrupole pure rotational transition \\
\textbf{Primary tracers} & warm molecular hydrogen mass; low-temperature rotational excitation; ortho-to-para ratio with other lines \\
\textbf{Typical environments} & warm molecular clouds; PDRs; shocks; galaxy disks \\
\textbf{Indicative regime} & sensitive to roughly 80-300 K gas; upper-state energy is about 510 K \\
\textbf{Density guidance} & Critical density depends on H, H2 and He collision rates, typically around 10\textasciicircum{}2-10\textasciicircum{}3 cm\textasciicircum{}-3 at warm temperatures. \\
\end{tabularx}
\end{conceptbox}

H2 0-0 S(0) is produced by electric-quadrupole pure rotational transition and is encountered in warm molecular clouds; PDRs; shocks; galaxy disks. It is commonly used to trace warm molecular hydrogen mass; low-temperature rotational excitation; ortho-to-para ratio with other lines.

\subsection*{How to use it}
Use this transition with its diagnostic companions to constrain warm molecular hydrogen mass; low-temperature rotational excitation; ortho-to-para ratio with other lines, while fitting the continuum, line-spread function, and relevant blends.

\subsection*{Interpretive cautions}
Quadrupole emission is optically thin but weak. Warm-H2 mass is exponentially sensitive to assumed temperature and ortho-to-para ratio.

\subsection*{Useful companions}
H2 0-0 S(1) 17.03484 micron; higher H2 rotational lines; CO rotational ladder

\subsection*{Sources}
\texttt{MPE\_H2; DRAINE\_2011; HOLLENBACH\_TIELENS\_1999}

\clearpage


\section{H2 0-0 S(1)}
\index{H2}

\begin{conceptbox}[At a glance]
\begin{tabularx}{\textwidth}{@{}lX@{}}
\textbf{Rest coordinate} & 17.03484 micron \\
\textbf{Band} & mid-infrared \\
\textbf{Transition} & v=0, J=3 -> 1 (ortho-H2) \\
\textbf{Mechanism} & electric-quadrupole pure rotational transition \\
\textbf{Primary tracers} & warm H2; shock/PDR excitation; rotational temperature \\
\textbf{Typical environments} & PDRs; molecular shocks; turbulent galaxy disks; AGN-heated molecular gas \\
\textbf{Indicative regime} & commonly traces about 150-700 K gas; upper-state energy is about 1015 K \\
\textbf{Density guidance} & Critical density is typically around 10\textasciicircum{}3-10\textasciicircum{}4 cm\textasciicircum{}-3, depending strongly on collider and temperature. \\
\end{tabularx}
\end{conceptbox}

H2 0-0 S(1) is produced by electric-quadrupole pure rotational transition and is encountered in PDRs; molecular shocks; turbulent galaxy disks; AGN-heated molecular gas. It is commonly used to trace warm H2; shock/PDR excitation; rotational temperature.

\subsection*{How to use it}
Use this transition with its diagnostic companions to constrain warm H2; shock/PDR excitation; rotational temperature, while fitting the continuum, line-spread function, and relevant blends.

\subsection*{Interpretive cautions}
Multiple-temperature excitation diagrams are usually required; a single Boltzmann slope can hide non-LTE or mixed gas.

\subsection*{Useful companions}
H2 0-0 S(0) 28.21883 micron; H2 0-0 S(2)/S(3); [Fe II] shock lines

\subsection*{Sources}
\texttt{MPE\_H2; DRAINE\_2011; HOLLENBACH\_TIELENS\_1999}

\clearpage


\section{H2 1-0 S(1)}
\index{H2}

\begin{conceptbox}[At a glance]
\begin{tabularx}{\textwidth}{@{}lX@{}}
\textbf{Rest coordinate} & 2.12183 micron \\
\textbf{Band} & near-infrared / K band \\
\textbf{Transition} & v=1, J=3 -> v=0, J=1 \\
\textbf{Mechanism} & rovibrational electric-quadrupole transition \\
\textbf{Primary tracers} & hot molecular gas; shock fronts; fluorescent versus thermal excitation with line ratios \\
\textbf{Typical environments} & molecular shocks; PDRs; protostellar jets; AGN and starburst nuclei \\
\textbf{Indicative regime} & roughly 1000-3000 K for thermal emission; UV fluorescence can populate it in cooler gas \\
\textbf{Density guidance} & Thermalization requires densities roughly 10\textasciicircum{}4-10\textasciicircum{}6 cm\textasciicircum{}-3 depending on collider and temperature; ratios carry the diagnostic power. \\
\end{tabularx}
\end{conceptbox}

H2 1-0 S(1) is produced by rovibrational electric-quadrupole transition and is encountered in molecular shocks; PDRs; protostellar jets; AGN and starburst nuclei. It is commonly used to trace hot molecular gas; shock fronts; fluorescent versus thermal excitation with line ratios.

\subsection*{How to use it}
Use this transition with its diagnostic companions to constrain hot molecular gas; shock fronts; fluorescent versus thermal excitation with line ratios, while fitting the continuum, line-spread function, and relevant blends.

\subsection*{Interpretive cautions}
2-1 S(1)/1-0 S(1) helps distinguish fluorescence from thermal shocks, but mixed excitation and differential extinction matter.

\subsection*{Useful companions}
H2 2-1 S(1) 2.2477 micron; Br-gamma 2.1661 micron; [Fe II] 1.644 micron

\subsection*{Sources}
\texttt{MPE\_H2; DRAINE\_2011; HOLLENBACH\_TIELENS\_1999}

\clearpage


\section{H I Brackett-gamma}
\index{H I}

\begin{conceptbox}[At a glance]
\begin{tabularx}{\textwidth}{@{}lX@{}}
\textbf{Rest coordinate} & 2.16612 micron \\
\textbf{Band} & near-infrared / K band \\
\textbf{Transition} & n=7 -> 4 \\
\textbf{Mechanism} & hydrogen recombination transition \\
\textbf{Primary tracers} & ionizing-photon rate; obscured star formation; accretion/outflow kinematics \\
\textbf{Typical environments} & dust-obscured H II regions; young stellar objects; stellar winds; AGN broad/narrow-line regions \\
\textbf{Indicative regime} & photoionized gas near 5000-20000 K \\
\textbf{Density guidance} & Case-B emissivity has modest density dependence over ordinary nebular conditions; high densities and optical depth alter decrements. \\
\end{tabularx}
\end{conceptbox}

H I Brackett-gamma is produced by hydrogen recombination transition and is encountered in dust-obscured H II regions; young stellar objects; stellar winds; AGN broad/narrow-line regions. It is commonly used to trace ionizing-photon rate; obscured star formation; accretion/outflow kinematics.

\subsection*{How to use it}
Use this transition with its diagnostic companions to constrain ionizing-photon rate; obscured star formation; accretion/outflow kinematics, while fitting the continuum, line-spread function, and relevant blends.

\subsection*{Interpretive cautions}
Much less extinguished than H-alpha but fainter. Correct for underlying stellar absorption, extinction and broad/narrow component mixing.

\subsection*{Useful companions}
Pa-alpha 1.8756 micron; H-alpha/H-beta; H2 2.1218 micron

\subsection*{Sources}
\texttt{NIST\_ASD; OSTERBROCK\_FERLAND\_2006; KENNICUTT\_EVANS\_2012}

\clearpage


\section{H I Paschen-alpha}
\index{H I}

\begin{conceptbox}[At a glance]
\begin{tabularx}{\textwidth}{@{}lX@{}}
\textbf{Rest coordinate} & 1.87561 micron \\
\textbf{Band} & near-infrared \\
\textbf{Transition} & n=4 -> 3 \\
\textbf{Mechanism} & hydrogen recombination transition \\
\textbf{Primary tracers} & ionizing-photon rate; dust-obscured star formation; extinction through recombination decrements \\
\textbf{Typical environments} & embedded H II regions; dusty starbursts; AGN line regions \\
\textbf{Indicative regime} & photoionized gas near 5000-20000 K \\
\textbf{Density guidance} & Weak density dependence in Case B; departures arise at high density, high optical depth or strong collisional excitation. \\
\end{tabularx}
\end{conceptbox}

H I Paschen-alpha is produced by hydrogen recombination transition and is encountered in embedded H II regions; dusty starbursts; AGN line regions. It is commonly used to trace ionizing-photon rate; dust-obscured star formation; extinction through recombination decrements.

\subsection*{How to use it}
Use this transition with its diagnostic companions to constrain ionizing-photon rate; dust-obscured star formation; extinction through recombination decrements, while fitting the continuum, line-spread function, and relevant blends.

\subsection*{Interpretive cautions}
Ground-based access is difficult because of telluric absorption; space observations are powerful for heavily obscured star formation.

\subsection*{Useful companions}
Br-gamma 2.1661 micron; H-alpha 6564.6 A; Pa-beta 1.2818 micron

\subsection*{Sources}
\texttt{NIST\_ASD; OSTERBROCK\_FERLAND\_2006; KENNICUTT\_EVANS\_2012}

\clearpage


\section{[Fe II] 1.644 micron}
\index{Fe II}

\begin{conceptbox}[At a glance]
\begin{tabularx}{\textwidth}{@{}lX@{}}
\textbf{Rest coordinate} & 1.644 micron \\
\textbf{Band} & near-infrared / H band \\
\textbf{Transition} & a4D7/2 -> a4F9/2 \\
\textbf{Mechanism} & forbidden fine-structure transition \\
\textbf{Primary tracers} & shocked/partially ionized gas; grain destruction; electron density with other [Fe II] lines \\
\textbf{Typical environments} & radiative shocks; supernova remnants; jets; partially ionized zones; AGN narrow-line regions \\
\textbf{Indicative regime} & typically several thousand to about 15000 K \\
\textbf{Density guidance} & Critical density is of order 10\textasciicircum{}4-10\textasciicircum{}5 cm\textasciicircum{}-3; ratios among lines from different upper levels constrain density. \\
\end{tabularx}
\end{conceptbox}

[Fe II] 1.644 micron is produced by forbidden fine-structure transition and is encountered in radiative shocks; supernova remnants; jets; partially ionized zones; AGN narrow-line regions. It is commonly used to trace shocked/partially ionized gas; grain destruction; electron density with other [Fe II] lines.

\subsection*{How to use it}
Use this transition with its diagnostic companions to constrain shocked/partially ionized gas; grain destruction; electron density with other [Fe II] lines, while fitting the continuum, line-spread function, and relevant blends.

\subsection*{Interpretive cautions}
The 1.257/1.644 ratio shares an upper level and is useful for extinction, but published A-values introduce systematic uncertainty.

\subsection*{Useful companions}
[Fe II] 1.257 micron; H2 2.1218 micron; Br-gamma

\subsection*{Sources}
\texttt{NIST\_ASD; OSTERBROCK\_FERLAND\_2006; DRAINE\_2011}

\clearpage


\section{He I 10830}
\index{He I}

\begin{conceptbox}[At a glance]
\begin{tabularx}{\textwidth}{@{}lX@{}}
\textbf{Rest coordinate} & 10833.3 angstrom \\
\textbf{Band} & near-infrared / Y-J boundary \\
\textbf{Transition} & 2s 3S -> 2p 3P multiplet \\
\textbf{Mechanism} & permitted transition from the metastable triplet state \\
\textbf{Primary tracers} & outflow velocity and covering fraction; metastable-helium column; hard EUV irradiation \\
\textbf{Typical environments} & stellar chromospheres; winds and accretion funnels; H II regions; AGN outflows; exoplanet escaping atmospheres \\
\textbf{Indicative regime} & photoionized/recombining gas from roughly 5000 K upward; metastable population is radiation and density dependent \\
\textbf{Density guidance} & Metastable 2s 3S population responds to electron density, ionization and radiative transfer; not a simple critical-density line. \\
\end{tabularx}
\end{conceptbox}

He I 10830 is produced by permitted transition from the metastable triplet state and is encountered in stellar chromospheres; winds and accretion funnels; H II regions; AGN outflows; exoplanet escaping atmospheres. It is commonly used to trace outflow velocity and covering fraction; metastable-helium column; hard EUV irradiation.

\subsection*{How to use it}
Use this transition with its diagnostic companions to constrain outflow velocity and covering fraction; metastable-helium column; hard EUV irradiation, while fitting the continuum, line-spread function, and relevant blends.

\subsection*{Interpretive cautions}
A multiplet with close components. Absorption can be saturated with partial covering; telluric correction and air/vacuum convention are essential.

\subsection*{Useful companions}
He I 5877.25 A; He I 2.058 micron; H-alpha; He II 4687 A

\subsection*{Sources}
\texttt{NIST\_ASD; OSTERBROCK\_FERLAND\_2006}

\clearpage


\section{[S III] 9533}
\index{S III}

\begin{conceptbox}[At a glance]
\begin{tabularx}{\textwidth}{@{}lX@{}}
\textbf{Rest coordinate} & 9533.2 angstrom \\
\textbf{Band} & optical / near-infrared \\
\textbf{Transition} & 1D2 -> 3P2 \\
\textbf{Mechanism} & forbidden collisionally excited transition \\
\textbf{Primary tracers} & S++ abundance; ionization/excitation; electron temperature with auroral [S III] 6312 \\
\textbf{Typical environments} & H II regions; planetary nebulae; star-forming galaxies \\
\textbf{Indicative regime} & photoionized gas around 5000-15000 K \\
\textbf{Density guidance} & Critical density is high enough that normal H II regions are usually in the low-density regime; IR fine-structure lines provide stronger density leverage. \\
\end{tabularx}
\end{conceptbox}

[S III] 9533 is produced by forbidden collisionally excited transition and is encountered in H II regions; planetary nebulae; star-forming galaxies. It is commonly used to trace S++ abundance; ionization/excitation; electron temperature with auroral [S III] 6312.

\subsection*{How to use it}
Use this transition with its diagnostic companions to constrain S++ abundance; ionization/excitation; electron temperature with auroral [S III] 6312, while fitting the continuum, line-spread function, and relevant blends.

\subsection*{Interpretive cautions}
Intrinsic 9533/9071 branching ratio is near 2.44 and can check calibration. Strong telluric absorption/fringing often dominates uncertainty.

\subsection*{Useful companions}
[S III] 9071 A; [S III] 6312 A; [S II] 6718/6733 A

\subsection*{Sources}
\texttt{NIST\_ASD; OSTERBROCK\_FERLAND\_2006}

\clearpage


\section{[S III] 9071}
\index{S III}

\begin{conceptbox}[At a glance]
\begin{tabularx}{\textwidth}{@{}lX@{}}
\textbf{Rest coordinate} & 9071.1 angstrom \\
\textbf{Band} & optical / near-infrared \\
\textbf{Transition} & 1D2 -> 3P1 \\
\textbf{Mechanism} & forbidden collisionally excited transition \\
\textbf{Primary tracers} & S++ abundance; ionized gas; electron temperature with [S III] 6312 \\
\textbf{Typical environments} & H II regions; planetary nebulae; star-forming galaxies \\
\textbf{Indicative regime} & photoionized gas around 5000-15000 K \\
\textbf{Density guidance} & Usually below its critical density in ordinary nebulae; ratio to 9533 is fixed mainly by A-values. \\
\end{tabularx}
\end{conceptbox}

[S III] 9071 is produced by forbidden collisionally excited transition and is encountered in H II regions; planetary nebulae; star-forming galaxies. It is commonly used to trace S++ abundance; ionized gas; electron temperature with [S III] 6312.

\subsection*{How to use it}
Use this transition with its diagnostic companions to constrain S++ abundance; ionized gas; electron temperature with [S III] 6312, while fitting the continuum, line-spread function, and relevant blends.

\subsection*{Interpretive cautions}
Vacuum value shown; the familiar air wavelength is about 9068.6 A. Correct for telluric absorption and use the theoretical branching ratio.

\subsection*{Useful companions}
[S III] 9533 A; [S III] 6312 A; [O II] 3727/3729 A

\subsection*{Sources}
\texttt{NIST\_ASD; OSTERBROCK\_FERLAND\_2006}

\clearpage


\section{H-alpha}
\index{H I}

\begin{conceptbox}[At a glance]
\begin{tabularx}{\textwidth}{@{}lX@{}}
\textbf{Rest coordinate} & 6564.61 angstrom \\
\textbf{Band} & optical / red \\
\textbf{Transition} & Balmer n=3 -> 2 \\
\textbf{Mechanism} & hydrogen recombination with possible collisional contribution \\
\textbf{Primary tracers} & recent massive-star formation; ionized-gas morphology and kinematics; recombination rate \\
\textbf{Typical environments} & H II regions; diffuse ionized gas; star-forming galaxies; AGN broad/narrow-line regions; radiative shocks \\
\textbf{Indicative regime} & photoionized gas near 5000-20000 K \\
\textbf{Density guidance} & Case-B emissivity scales approximately with n\_e n\_p but the H-alpha/H-beta decrement is mainly extinction/transfer sensitive over ordinary nebular densities. \\
\end{tabularx}
\end{conceptbox}

H-alpha is produced by hydrogen recombination with possible collisional contribution and is encountered in H II regions; diffuse ionized gas; star-forming galaxies; AGN broad/narrow-line regions; radiative shocks. It is commonly used to trace recent massive-star formation; ionized-gas morphology and kinematics; recombination rate.

\subsection*{How to use it}
Use this transition with its diagnostic companions to constrain recent massive-star formation; ionized-gas morphology and kinematics; recombination rate, while fitting the continuum, line-spread function, and relevant blends.

\subsection*{Interpretive cautions}
Vacuum value shown; the traditional air wavelength is 6562.8 A. Dust, [N II] blending, stellar absorption, leakage and stochastic IMF sampling affect SFR estimates.

\subsection*{Useful companions}
H-beta 4862.68 A for Balmer decrement; [N II] 6585 A and [S II] 6718/6733 A for BPT/shock diagnostics; Pa-alpha/Br-gamma

\subsection*{Sources}
\texttt{NIST\_ASD; KENNICUTT\_EVANS\_2012; OSTERBROCK\_FERLAND\_2006}

\clearpage


\section{[N II] 6585}
\index{N II}

\begin{conceptbox}[At a glance]
\begin{tabularx}{\textwidth}{@{}lX@{}}
\textbf{Rest coordinate} & 6585.27 angstrom \\
\textbf{Band} & optical / red \\
\textbf{Transition} & 1D2 -> 3P2 \\
\textbf{Mechanism} & forbidden collisionally excited transition \\
\textbf{Primary tracers} & gas-phase metallicity through calibrated ratios; excitation source with H-alpha; N+ abundance \\
\textbf{Typical environments} & H II regions; diffuse ionized gas; planetary nebulae; AGN narrow-line regions; shocks \\
\textbf{Indicative regime} & photoionized or shocked gas around 6000-20000 K \\
\textbf{Density guidance} & Critical density is roughly 10\textasciicircum{}5 cm\textasciicircum{}-3; usually in the low-density regime for normal H II regions. \\
\end{tabularx}
\end{conceptbox}

[N II] 6585 is produced by forbidden collisionally excited transition and is encountered in H II regions; diffuse ionized gas; planetary nebulae; AGN narrow-line regions; shocks. It is commonly used to trace gas-phase metallicity through calibrated ratios; excitation source with H-alpha; N+ abundance.

\subsection*{How to use it}
Use this transition with its diagnostic companions to constrain gas-phase metallicity through calibrated ratios; excitation source with H-alpha; N+ abundance, while fitting the continuum, line-spread function, and relevant blends.

\subsection*{Interpretive cautions}
Vacuum value; familiar air wavelength is 6583.45 A. N/O variations and ionization parameter limit one-dimensional metallicity calibrations.

\subsection*{Useful companions}
[N II] 6549.86 A fixed branching partner; H-alpha for N2/BPT; [N II] 5756 A for electron temperature

\subsection*{Sources}
\texttt{NIST\_ASD; KEWLEY\_2006; OSTERBROCK\_FERLAND\_2006}

\clearpage


\section{[S II] 6718}
\index{S II}

\begin{conceptbox}[At a glance]
\begin{tabularx}{\textwidth}{@{}lX@{}}
\textbf{Rest coordinate} & 6718.29 angstrom \\
\textbf{Band} & optical / red \\
\textbf{Transition} & 2D5/2 -> 4S3/2 \\
\textbf{Mechanism} & forbidden collisionally excited transition \\
\textbf{Primary tracers} & electron density with 6733; low-ionization/shock zones; excitation source with H-alpha \\
\textbf{Typical environments} & H II regions; diffuse ionized gas; shock fronts; AGN narrow-line regions; supernova remnants \\
\textbf{Indicative regime} & roughly 5000-20000 K ionized or partially ionized gas \\
\textbf{Density guidance} & Critical density near 10\textasciicircum{}4 K is about 1.5x10\textasciicircum{}3 cm\textasciicircum{}-3; the doublet ratio is most useful around 10\textasciicircum{}2-10\textasciicircum{}4 cm\textasciicircum{}-3. \\
\end{tabularx}
\end{conceptbox}

[S II] 6718 is produced by forbidden collisionally excited transition and is encountered in H II regions; diffuse ionized gas; shock fronts; AGN narrow-line regions; supernova remnants. It is commonly used to trace electron density with 6733; low-ionization/shock zones; excitation source with H-alpha.

\subsection*{How to use it}
Use this transition with its diagnostic companions to constrain electron density with 6733; low-ionization/shock zones; excitation source with H-alpha, while fitting the continuum, line-spread function, and relevant blends.

\subsection*{Interpretive cautions}
Vacuum value; familiar air wavelength is 6716.44 A. Ratio saturates at low and high density and can mix unresolved density phases.

\subsection*{Useful companions}
[S II] 6733 A; H-alpha for BPT/shock ratios; [O I] 6302 A

\subsection*{Sources}
\texttt{NIST\_ASD; KEWLEY\_2006; OSTERBROCK\_FERLAND\_2006}

\clearpage


\section{[S II] 6733}
\index{S II}

\begin{conceptbox}[At a glance]
\begin{tabularx}{\textwidth}{@{}lX@{}}
\textbf{Rest coordinate} & 6732.67 angstrom \\
\textbf{Band} & optical / red \\
\textbf{Transition} & 2D3/2 -> 4S3/2 \\
\textbf{Mechanism} & forbidden collisionally excited transition \\
\textbf{Primary tracers} & electron density with 6718; low-ionization gas; shock excitation \\
\textbf{Typical environments} & H II regions; diffuse ionized gas; shock fronts; AGN narrow-line regions; supernova remnants \\
\textbf{Indicative regime} & roughly 5000-20000 K \\
\textbf{Density guidance} & Critical density near 10\textasciicircum{}4 K is about 4x10\textasciicircum{}3 cm\textasciicircum{}-3; paired ratio carries the density information. \\
\end{tabularx}
\end{conceptbox}

[S II] 6733 is produced by forbidden collisionally excited transition and is encountered in H II regions; diffuse ionized gas; shock fronts; AGN narrow-line regions; supernova remnants. It is commonly used to trace electron density with 6718; low-ionization gas; shock excitation.

\subsection*{How to use it}
Use this transition with its diagnostic companions to constrain electron density with 6718; low-ionization gas; shock excitation, while fitting the continuum, line-spread function, and relevant blends.

\subsection*{Interpretive cautions}
Vacuum value; familiar air wavelength is 6730.82 A. Deblend both lines from sky residuals and broad H-alpha wings in complex nuclei.

\subsection*{Useful companions}
[S II] 6718 A; H-alpha; [O II] 3727/3729 A

\subsection*{Sources}
\texttt{NIST\_ASD; KEWLEY\_2006; OSTERBROCK\_FERLAND\_2006}

\clearpage


\section{[O I] 6302}
\index{O I}

\begin{conceptbox}[At a glance]
\begin{tabularx}{\textwidth}{@{}lX@{}}
\textbf{Rest coordinate} & 6302.05 angstrom \\
\textbf{Band} & optical / red \\
\textbf{Transition} & 1D2 -> 3P2 \\
\textbf{Mechanism} & forbidden collisionally excited transition \\
\textbf{Primary tracers} & hard ionizing spectra; shock excitation; neutral/ionized interface \\
\textbf{Typical environments} & partially ionized zones; shocks; AGN narrow-line regions; nebular ionization fronts; night-sky airglow \\
\textbf{Indicative regime} & warm neutral or partially ionized gas around several thousand to 20000 K \\
\textbf{Density guidance} & Electron critical density is of order 2x10\textasciicircum{}6 cm\textasciicircum{}-3; collisional quenching matters in dense gas. \\
\end{tabularx}
\end{conceptbox}

[O I] 6302 is produced by forbidden collisionally excited transition and is encountered in partially ionized zones; shocks; AGN narrow-line regions; nebular ionization fronts; night-sky airglow. It is commonly used to trace hard ionizing spectra; shock excitation; neutral/ionized interface.

\subsection*{How to use it}
Use this transition with its diagnostic companions to constrain hard ionizing spectra; shock excitation; neutral/ionized interface, while fitting the continuum, line-spread function, and relevant blends.

\subsection*{Interpretive cautions}
Vacuum value; familiar air wavelength is 6300.30 A. Bright terrestrial airglow demands careful sky subtraction, especially at low redshift.

\subsection*{Useful companions}
[O I] 6365.54 A fixed branching partner; H-alpha for BPT; [S II] doublet

\subsection*{Sources}
\texttt{NIST\_ASD; KEWLEY\_2006; OSTERBROCK\_FERLAND\_2006}

\clearpage


\section{He I 5877}
\index{He I}

\begin{conceptbox}[At a glance]
\begin{tabularx}{\textwidth}{@{}lX@{}}
\textbf{Rest coordinate} & 5877.25 angstrom \\
\textbf{Band} & optical / yellow \\
\textbf{Transition} & 3d 3D -> 2p 3P multiplet \\
\textbf{Mechanism} & helium recombination transition with collisional and radiative-transfer corrections \\
\textbf{Primary tracers} & He+/H+ abundance; hardness above 24.6 eV; nebular physical conditions with other He I lines \\
\textbf{Typical environments} & H II regions; planetary nebulae; hot-star winds; AGN line regions \\
\textbf{Indicative regime} & photoionized gas around 5000-20000 K \\
\textbf{Density guidance} & He I line ratios respond to electron density, optical depth in triplet lines and underlying absorption; multi-line fitting is required. \\
\end{tabularx}
\end{conceptbox}

He I 5877 is produced by helium recombination transition with collisional and radiative-transfer corrections and is encountered in H II regions; planetary nebulae; hot-star winds; AGN line regions. It is commonly used to trace He+/H+ abundance; hardness above 24.6 eV; nebular physical conditions with other He I lines.

\subsection*{How to use it}
Use this transition with its diagnostic companions to constrain He+/H+ abundance; hardness above 24.6 eV; nebular physical conditions with other He I lines, while fitting the continuum, line-spread function, and relevant blends.

\subsection*{Interpretive cautions}
Vacuum value; familiar air wavelength is about 5875.62 A. Correct for stellar absorption, Na sky features and triplet optical-depth effects.

\subsection*{Useful companions}
He I 6679 A; He I 7067 A; He I 10833 A; H-beta

\subsection*{Sources}
\texttt{NIST\_ASD; OSTERBROCK\_FERLAND\_2006}

\clearpage


\section{[O III] 5008}
\index{O III}

\begin{conceptbox}[At a glance]
\begin{tabularx}{\textwidth}{@{}lX@{}}
\textbf{Rest coordinate} & 5008.24 angstrom \\
\textbf{Band} & optical / green \\
\textbf{Transition} & 1D2 -> 3P2 \\
\textbf{Mechanism} & forbidden collisionally excited transition \\
\textbf{Primary tracers} & high-excitation ionized gas; AGN power statistically; metallicity/ionization with other lines; kinematics \\
\textbf{Typical environments} & H II regions; planetary nebulae; AGN narrow-line regions; ionized outflows \\
\textbf{Indicative regime} & photoionized gas about 7000-20000 K; shocks can extend higher \\
\textbf{Density guidance} & Critical density is about 7x10\textasciicircum{}5 cm\textasciicircum{}-3; the nebular-to-auroral ratio with 4364 A is strongly temperature sensitive. \\
\end{tabularx}
\end{conceptbox}

[O III] 5008 is produced by forbidden collisionally excited transition and is encountered in H II regions; planetary nebulae; AGN narrow-line regions; ionized outflows. It is commonly used to trace high-excitation ionized gas; AGN power statistically; metallicity/ionization with other lines; kinematics.

\subsection*{How to use it}
Use this transition with its diagnostic companions to constrain high-excitation ionized gas; AGN power statistically; metallicity/ionization with other lines; kinematics, while fitting the continuum, line-spread function, and relevant blends.

\subsection*{Interpretive cautions}
Vacuum value; familiar air wavelength is 5006.84 A. Use the theoretical 5008/4960 ratio near 2.98 to test fits and saturation.

\subsection*{Useful companions}
[O III] 4960 A fixed branching partner; [O III] 4364 A temperature line; H-beta for BPT/R3; [O II] for O32/R23

\subsection*{Sources}
\texttt{NIST\_ASD; KEWLEY\_2006; OSTERBROCK\_FERLAND\_2006}

\clearpage


\section{[O III] 4960}
\index{O III}

\begin{conceptbox}[At a glance]
\begin{tabularx}{\textwidth}{@{}lX@{}}
\textbf{Rest coordinate} & 4960.3 angstrom \\
\textbf{Band} & optical / blue-green \\
\textbf{Transition} & 1D2 -> 3P1 \\
\textbf{Mechanism} & forbidden collisionally excited transition \\
\textbf{Primary tracers} & O++ abundance; high-excitation gas; quality control with 5008 \\
\textbf{Typical environments} & H II regions; planetary nebulae; AGN narrow-line regions \\
\textbf{Indicative regime} & photoionized gas about 7000-20000 K \\
\textbf{Density guidance} & Same upper level and critical-density behavior as 5008 A. \\
\end{tabularx}
\end{conceptbox}

[O III] 4960 is produced by forbidden collisionally excited transition and is encountered in H II regions; planetary nebulae; AGN narrow-line regions. It is commonly used to trace O++ abundance; high-excitation gas; quality control with 5008.

\subsection*{How to use it}
Use this transition with its diagnostic companions to constrain O++ abundance; high-excitation gas; quality control with 5008, while fitting the continuum, line-spread function, and relevant blends.

\subsection*{Interpretive cautions}
Vacuum value; familiar air wavelength is 4958.91 A. Tie centroid, width and theoretical flux ratio to 5008 in multi-component fits when justified.

\subsection*{Useful companions}
[O III] 5008 A; [O III] 4364 A; H-beta

\subsection*{Sources}
\texttt{NIST\_ASD; OSTERBROCK\_FERLAND\_2006}

\clearpage


\section{H-beta}
\index{H I}

\begin{conceptbox}[At a glance]
\begin{tabularx}{\textwidth}{@{}lX@{}}
\textbf{Rest coordinate} & 4862.68 angstrom \\
\textbf{Band} & optical / blue \\
\textbf{Transition} & Balmer n=4 -> 2 \\
\textbf{Mechanism} & hydrogen recombination with possible collisional contribution \\
\textbf{Primary tracers} & recombination rate; dust reddening with H-alpha; normalization for excitation and abundance ratios \\
\textbf{Typical environments} & H II regions; star-forming galaxies; planetary nebulae; AGN line regions \\
\textbf{Indicative regime} & photoionized gas near 5000-20000 K \\
\textbf{Density guidance} & Case-B ratios depend weakly on density over normal nebular conditions but change with temperature, optical depth and collisions. \\
\end{tabularx}
\end{conceptbox}

H-beta is produced by hydrogen recombination with possible collisional contribution and is encountered in H II regions; star-forming galaxies; planetary nebulae; AGN line regions. It is commonly used to trace recombination rate; dust reddening with H-alpha; normalization for excitation and abundance ratios.

\subsection*{How to use it}
Use this transition with its diagnostic companions to constrain recombination rate; dust reddening with H-alpha; normalization for excitation and abundance ratios, while fitting the continuum, line-spread function, and relevant blends.

\subsection*{Interpretive cautions}
Vacuum value; familiar air wavelength is 4861.33 A. Stellar Balmer absorption can severely bias weak emission and extinction estimates.

\subsection*{Useful companions}
H-alpha Balmer decrement; [O III] 5008/H-beta; H-gamma absorption/emission correction

\subsection*{Sources}
\texttt{NIST\_ASD; KEWLEY\_2006; OSTERBROCK\_FERLAND\_2006}

\clearpage


\section{He II 4687}
\index{He II}

\begin{conceptbox}[At a glance]
\begin{tabularx}{\textwidth}{@{}lX@{}}
\textbf{Rest coordinate} & 4687.02 angstrom \\
\textbf{Band} & optical / blue \\
\textbf{Transition} & hydrogenic n=4 -> 3 \\
\textbf{Mechanism} & recombination transition of He+ \\
\textbf{Primary tracers} & photons above 54.4 eV; very hard ionizing spectra; He++ abundance \\
\textbf{Typical environments} & planetary nebulae; Wolf-Rayet populations; X-ray binaries; AGN; very low-metallicity star-forming regions; fast shocks \\
\textbf{Indicative regime} & photoionized gas commonly 10000-30000 K; shock zones can be hotter \\
\textbf{Density guidance} & Recombination emissivity rather than a critical-density diagnostic; line width separates nebular and stellar components. \\
\end{tabularx}
\end{conceptbox}

He II 4687 is produced by recombination transition of He+ and is encountered in planetary nebulae; Wolf-Rayet populations; X-ray binaries; AGN; very low-metallicity star-forming regions; fast shocks. It is commonly used to trace photons above 54.4 eV; very hard ionizing spectra; He++ abundance.

\subsection*{How to use it}
Use this transition with its diagnostic companions to constrain photons above 54.4 eV; very hard ionizing spectra; He++ abundance, while fitting the continuum, line-spread function, and relevant blends.

\subsection*{Interpretive cautions}
Vacuum value; familiar air wavelength is 4685.7 A. Broad stellar and narrow nebular He II require decomposition before inferring ionizing hardness.

\subsection*{Useful companions}
He II 1640 A; H-beta; [Ne V] or high-ionization lines; Wolf-Rayet blue bump

\subsection*{Sources}
\texttt{NIST\_ASD; OSTERBROCK\_FERLAND\_2006}

\clearpage


\section{[O II] 3727}
\index{O II}

\begin{conceptbox}[At a glance]
\begin{tabularx}{\textwidth}{@{}lX@{}}
\textbf{Rest coordinate} & 3727.09 angstrom \\
\textbf{Band} & optical / near-UV \\
\textbf{Transition} & 2D3/2 -> 4S3/2 (traditional 3726 component) \\
\textbf{Mechanism} & forbidden collisionally excited transition \\
\textbf{Primary tracers} & electron density with 3729; O+ abundance; star formation after calibration; ionization parameter with [O III] \\
\textbf{Typical environments} & H II regions; star-forming galaxies; diffuse ionized gas; AGN narrow-line regions \\
\textbf{Indicative regime} & photoionized gas roughly 7000-20000 K \\
\textbf{Density guidance} & Critical density is roughly 1.4x10\textasciicircum{}4 cm\textasciicircum{}-3 near 10\textasciicircum{}4 K; doublet ratio spans common nebular densities. \\
\end{tabularx}
\end{conceptbox}

[O II] 3727 is produced by forbidden collisionally excited transition and is encountered in H II regions; star-forming galaxies; diffuse ionized gas; AGN narrow-line regions. It is commonly used to trace electron density with 3729; O+ abundance; star formation after calibration; ionization parameter with [O III].

\subsection*{How to use it}
Use this transition with its diagnostic companions to constrain electron density with 3729; O+ abundance; star formation after calibration; ionization parameter with [O III], while fitting the continuum, line-spread function, and relevant blends.

\subsection*{Interpretive cautions}
Vacuum value; traditional air wavelength is 3726.03 A. Unresolved surveys often quote [O II] 3727 for the sum—record that convention explicitly.

\subsection*{Useful companions}
[O II] 3729.88 A; [O III] 5008 A for O32/R23; [O II] 7320/7330 A temperature lines

\subsection*{Sources}
\texttt{NIST\_ASD; KENNICUTT\_EVANS\_2012; OSTERBROCK\_FERLAND\_2006}

\clearpage


\section{[O II] 3729}
\index{O II}

\begin{conceptbox}[At a glance]
\begin{tabularx}{\textwidth}{@{}lX@{}}
\textbf{Rest coordinate} & 3729.88 angstrom \\
\textbf{Band} & optical / near-UV \\
\textbf{Transition} & 2D5/2 -> 4S3/2 (traditional 3729 component) \\
\textbf{Mechanism} & forbidden collisionally excited transition \\
\textbf{Primary tracers} & electron density with 3727; O+ abundance; low-ionization gas \\
\textbf{Typical environments} & H II regions; star-forming galaxies; diffuse ionized gas; AGN narrow-line regions \\
\textbf{Indicative regime} & photoionized gas roughly 7000-20000 K \\
\textbf{Density guidance} & Critical density is roughly 3x10\textasciicircum{}3 cm\textasciicircum{}-3 near 10\textasciicircum{}4 K; ratio to 3727 saturates outside its useful density range. \\
\end{tabularx}
\end{conceptbox}

[O II] 3729 is produced by forbidden collisionally excited transition and is encountered in H II regions; star-forming galaxies; diffuse ionized gas; AGN narrow-line regions. It is commonly used to trace electron density with 3727; O+ abundance; low-ionization gas.

\subsection*{How to use it}
Use this transition with its diagnostic companions to constrain electron density with 3727; O+ abundance; low-ionization gas, while fitting the continuum, line-spread function, and relevant blends.

\subsection*{Interpretive cautions}
Vacuum value; traditional air wavelength is 3728.82 A. Instrumental resolution and stellar continuum modeling determine whether the doublet yields useful densities.

\subsection*{Useful companions}
[O II] 3727.09 A; [O III] 5008 A; [S II] 6718/6733 A

\subsection*{Sources}
\texttt{NIST\_ASD; OSTERBROCK\_FERLAND\_2006}

\clearpage


\section{[Ne III] 3869}
\index{Ne III}

\begin{conceptbox}[At a glance]
\begin{tabularx}{\textwidth}{@{}lX@{}}
\textbf{Rest coordinate} & 3869.86 angstrom \\
\textbf{Band} & optical / violet \\
\textbf{Transition} & 1D2 -> 3P2 \\
\textbf{Mechanism} & forbidden collisionally excited transition \\
\textbf{Primary tracers} & hard ionizing spectrum; Ne++ abundance; ionization parameter with [O II] or [Ne II] \\
\textbf{Typical environments} & high-excitation H II regions; planetary nebulae; AGN narrow-line regions \\
\textbf{Indicative regime} & photoionized gas around 8000-20000 K \\
\textbf{Density guidance} & Critical density is high, of order 10\textasciicircum{}7 cm\textasciicircum{}-3; generally in the low-density regime for H II regions. \\
\end{tabularx}
\end{conceptbox}

[Ne III] 3869 is produced by forbidden collisionally excited transition and is encountered in high-excitation H II regions; planetary nebulae; AGN narrow-line regions. It is commonly used to trace hard ionizing spectrum; Ne++ abundance; ionization parameter with [O II] or [Ne II].

\subsection*{How to use it}
Use this transition with its diagnostic companions to constrain hard ionizing spectrum; Ne++ abundance; ionization parameter with [O II] or [Ne II], while fitting the continuum, line-spread function, and relevant blends.

\subsection*{Interpretive cautions}
Vacuum value; familiar air wavelength is 3868.76 A. The weaker 3968-A partner blends with H-epsilon/Ca H in many spectra.

\subsection*{Useful companions}
[Ne III] 3968 A branching partner; [O II] 3727+3729; [Ne III] 15.6 micron

\subsection*{Sources}
\texttt{NIST\_ASD; OSTERBROCK\_FERLAND\_2006}

\clearpage


\section{Mg II 2796}
\index{Mg II}

\begin{conceptbox}[At a glance]
\begin{tabularx}{\textwidth}{@{}lX@{}}
\textbf{Rest coordinate} & 2796.352 angstrom \\
\textbf{Band} & near-ultraviolet \\
\textbf{Transition} & 3s 2S1/2 -> 3p 2P3/2 \\
\textbf{Mechanism} & permitted resonance transition \\
\textbf{Primary tracers} & cool photoionized gas; kinematics and covering fraction; outflows/inflows; magnesium column with saturation modeling \\
\textbf{Typical environments} & cool CGM; intervening absorbers; galactic winds; stellar chromospheres; quasar broad-line regions \\
\textbf{Indicative regime} & commonly about 10\textasciicircum{}4 K photoionized gas; not a unique thermometer \\
\textbf{Density guidance} & Resonance absorption is column/ionization sensitive, not a critical-density diagnostic; emission is shaped by resonant scattering. \\
\end{tabularx}
\end{conceptbox}

Mg II 2796 is produced by permitted resonance transition and is encountered in cool CGM; intervening absorbers; galactic winds; stellar chromospheres; quasar broad-line regions. It is commonly used to trace cool photoionized gas; kinematics and covering fraction; outflows/inflows; magnesium column with saturation modeling.

\subsection*{How to use it}
Use this transition with its diagnostic companions to constrain cool photoionized gas; kinematics and covering fraction; outflows/inflows; magnesium column with saturation modeling, while fitting the continuum, line-spread function, and relevant blends.

\subsection*{Interpretive cautions}
Stronger doublet member with oscillator-strength ratio near 2. Saturation and partial covering drive the observed doublet ratio toward unity.

\subsection*{Useful companions}
Mg II 2803.530 A; Mg I 2852.96 A; Fe II UV lines

\subsection*{Sources}
\texttt{NIST\_ASD; MORTON\_2003; TUMLINSON\_2017}

\clearpage


\section{Mg II 2803}
\index{Mg II}

\begin{conceptbox}[At a glance]
\begin{tabularx}{\textwidth}{@{}lX@{}}
\textbf{Rest coordinate} & 2803.53 angstrom \\
\textbf{Band} & near-ultraviolet \\
\textbf{Transition} & 3s 2S1/2 -> 3p 2P1/2 \\
\textbf{Mechanism} & permitted resonance transition \\
\textbf{Primary tracers} & cool gas column and kinematics; saturation/covering with 2796; outflows and inflows \\
\textbf{Typical environments} & cool CGM; intervening absorbers; galactic winds; stellar chromospheres; quasar broad-line regions \\
\textbf{Indicative regime} & commonly about 10\textasciicircum{}4 K photoionized gas \\
\textbf{Density guidance} & Not a direct density diagnostic; its value comes from doublet curve-of-growth and covering analysis. \\
\end{tabularx}
\end{conceptbox}

Mg II 2803 is produced by permitted resonance transition and is encountered in cool CGM; intervening absorbers; galactic winds; stellar chromospheres; quasar broad-line regions. It is commonly used to trace cool gas column and kinematics; saturation/covering with 2796; outflows and inflows.

\subsection*{How to use it}
Use this transition with its diagnostic companions to constrain cool gas column and kinematics; saturation/covering with 2796; outflows and inflows, while fitting the continuum, line-spread function, and relevant blends.

\subsection*{Interpretive cautions}
Weaker doublet member is essential for diagnosing saturation. Resonant emission infill can bias absorption equivalent widths.

\subsection*{Useful companions}
Mg II 2796.352 A; Mg I 2852.96 A; Fe II 2600.17 A

\subsection*{Sources}
\texttt{NIST\_ASD; MORTON\_2003; TUMLINSON\_2017}

\clearpage


\section{C III] 1907}
\index{C III}

\begin{conceptbox}[At a glance]
\begin{tabularx}{\textwidth}{@{}lX@{}}
\textbf{Rest coordinate} & 1906.683 angstrom \\
\textbf{Band} & ultraviolet \\
\textbf{Transition} & 2s2p 3P2 -> 2s2 1S0 \\
\textbf{Mechanism} & semi-forbidden intercombination transition \\
\textbf{Primary tracers} & electron density with 1909; carbon abundance and ionization; hard young stellar populations \\
\textbf{Typical environments} & H II regions; high-redshift star-forming galaxies; AGN narrow/broad-line regions; planetary nebulae \\
\textbf{Indicative regime} & photoionized gas around 10000-30000 K \\
\textbf{Density guidance} & The 1907/1909 ratio is density sensitive from low nebular densities into roughly 10\textasciicircum{}5-10\textasciicircum{}6 cm\textasciicircum{}-3, with isotope/atomic-data nuances. \\
\end{tabularx}
\end{conceptbox}

C III] 1907 is produced by semi-forbidden intercombination transition and is encountered in H II regions; high-redshift star-forming galaxies; AGN narrow/broad-line regions; planetary nebulae. It is commonly used to trace electron density with 1909; carbon abundance and ionization; hard young stellar populations.

\subsection*{How to use it}
Use this transition with its diagnostic companions to constrain electron density with 1909; carbon abundance and ionization; hard young stellar populations, while fitting the continuum, line-spread function, and relevant blends.

\subsection*{Interpretive cautions}
Usually unresolved at modest resolution and reported as C III] 1908. AGN broad-line density can suppress or alter the doublet.

\subsection*{Useful companions}
C III] 1908.734 A; O III] 1661/1666 A; He II 1640 A; C IV 1548/1551 A

\subsection*{Sources}
\texttt{NIST\_ASD; MORTON\_2003; OSTERBROCK\_FERLAND\_2006}

\clearpage


\section{C III] 1909}
\index{C III}

\begin{conceptbox}[At a glance]
\begin{tabularx}{\textwidth}{@{}lX@{}}
\textbf{Rest coordinate} & 1908.734 angstrom \\
\textbf{Band} & ultraviolet \\
\textbf{Transition} & 2s2p 3P1 -> 2s2 1S0 \\
\textbf{Mechanism} & semi-forbidden intercombination transition \\
\textbf{Primary tracers} & electron density; carbon abundance; ionizing conditions \\
\textbf{Typical environments} & H II regions; high-redshift galaxies; AGN line regions; planetary nebulae \\
\textbf{Indicative regime} & photoionized gas around 10000-30000 K \\
\textbf{Density guidance} & Pair with 1907; the ratio approaches atomic low-density and high-density limits and must not be extrapolated beyond them. \\
\end{tabularx}
\end{conceptbox}

C III] 1909 is produced by semi-forbidden intercombination transition and is encountered in H II regions; high-redshift galaxies; AGN line regions; planetary nebulae. It is commonly used to trace electron density; carbon abundance; ionizing conditions.

\subsection*{How to use it}
Use this transition with its diagnostic companions to constrain electron density; carbon abundance; ionizing conditions, while fitting the continuum, line-spread function, and relevant blends.

\subsection*{Interpretive cautions}
Blends with 1907 at common survey resolution; density, C/O, ionization parameter and stellar/AGN continuum are degenerate in equivalent width.

\subsection*{Useful companions}
C III] 1906.683 A; Si III] 1883/1892 A; O III] 1666 A

\subsection*{Sources}
\texttt{NIST\_ASD; MORTON\_2003; OSTERBROCK\_FERLAND\_2006}

\clearpage


\section{C IV 1548}
\index{C IV}

\begin{conceptbox}[At a glance]
\begin{tabularx}{\textwidth}{@{}lX@{}}
\textbf{Rest coordinate} & 1548.204 angstrom \\
\textbf{Band} & far-ultraviolet \\
\textbf{Transition} & 2s 2S1/2 -> 2p 2P3/2 \\
\textbf{Mechanism} & permitted resonance transition \\
\textbf{Primary tracers} & high-ionization gas; metal-enriched CGM; outflow kinematics; AGN/stars when in emission \\
\textbf{Typical environments} & warm CGM/IGM; quasar absorbers; AGN broad-line regions; hot-star winds; transition-temperature plasma \\
\textbf{Indicative regime} & CIE abundance peaks near 10\textasciicircum{}5 K; photoionized C IV commonly occurs near 10\textasciicircum{}4 K \\
\textbf{Density guidance} & Absorption constrains ion column and ionization, not density alone; emission can be resonantly scattered and collisionally excited. \\
\end{tabularx}
\end{conceptbox}

C IV 1548 is produced by permitted resonance transition and is encountered in warm CGM/IGM; quasar absorbers; AGN broad-line regions; hot-star winds; transition-temperature plasma. It is commonly used to trace high-ionization gas; metal-enriched CGM; outflow kinematics; AGN/stars when in emission.

\subsection*{How to use it}
Use this transition with its diagnostic companions to constrain high-ionization gas; metal-enriched CGM; outflow kinematics; AGN/stars when in emission, while fitting the continuum, line-spread function, and relevant blends.

\subsection*{Interpretive cautions}
Stronger doublet member. Saturation, unresolved components and partial covering are major column-density systematics; stellar P-Cygni and nebular/AGN emission can blend.

\subsection*{Useful companions}
C IV 1550.781 A; Si IV 1394/1403 A; N V 1239/1243 A; O VI 1032/1038 A

\subsection*{Sources}
\texttt{NIST\_ASD; MORTON\_2003; TUMLINSON\_2017}

\clearpage


\section{C IV 1551}
\index{C IV}

\begin{conceptbox}[At a glance]
\begin{tabularx}{\textwidth}{@{}lX@{}}
\textbf{Rest coordinate} & 1550.781 angstrom \\
\textbf{Band} & far-ultraviolet \\
\textbf{Transition} & 2s 2S1/2 -> 2p 2P1/2 \\
\textbf{Mechanism} & permitted resonance transition \\
\textbf{Primary tracers} & doublet saturation and covering; high-ionization gas kinematics; metal enrichment \\
\textbf{Typical environments} & warm CGM/IGM; quasar absorbers; AGN; hot-star winds \\
\textbf{Indicative regime} & CIE peak near 10\textasciicircum{}5 K; also common in cooler photoionized gas \\
\textbf{Density guidance} & No simple critical-density diagnostic in absorption; line ratio tests saturation and radiative transfer. \\
\end{tabularx}
\end{conceptbox}

C IV 1551 is produced by permitted resonance transition and is encountered in warm CGM/IGM; quasar absorbers; AGN; hot-star winds. It is commonly used to trace doublet saturation and covering; high-ionization gas kinematics; metal enrichment.

\subsection*{How to use it}
Use this transition with its diagnostic companions to constrain doublet saturation and covering; high-ionization gas kinematics; metal enrichment, while fitting the continuum, line-spread function, and relevant blends.

\subsection*{Interpretive cautions}
Weaker member has about half the oscillator strength of 1548. Intrinsic AGN doublet blending depends on line width.

\subsection*{Useful companions}
C IV 1548.204 A; Si IV doublet; O VI doublet

\subsection*{Sources}
\texttt{NIST\_ASD; MORTON\_2003; TUMLINSON\_2017}

\clearpage


\section{Si IV 1394}
\index{Si IV}

\begin{conceptbox}[At a glance]
\begin{tabularx}{\textwidth}{@{}lX@{}}
\textbf{Rest coordinate} & 1393.755 angstrom \\
\textbf{Band} & far-ultraviolet \\
\textbf{Transition} & 3s 2S1/2 -> 3p 2P3/2 \\
\textbf{Mechanism} & permitted resonance transition \\
\textbf{Primary tracers} & intermediate/high-ionization gas; outflow kinematics; silicon enrichment; stellar-wind profiles \\
\textbf{Typical environments} & photoionized CGM; galactic winds; IGM absorbers; hot-star winds; transition-temperature interfaces \\
\textbf{Indicative regime} & CIE abundance peaks around 6-8x10\textasciicircum{}4 K; photoionized gas can be near 10\textasciicircum{}4 K \\
\textbf{Density guidance} & Primarily column/ionization sensitive; the doublet tests saturation and covering fraction. \\
\end{tabularx}
\end{conceptbox}

Si IV 1394 is produced by permitted resonance transition and is encountered in photoionized CGM; galactic winds; IGM absorbers; hot-star winds; transition-temperature interfaces. It is commonly used to trace intermediate/high-ionization gas; outflow kinematics; silicon enrichment; stellar-wind profiles.

\subsection*{How to use it}
Use this transition with its diagnostic companions to constrain intermediate/high-ionization gas; outflow kinematics; silicon enrichment; stellar-wind profiles, while fitting the continuum, line-spread function, and relevant blends.

\subsection*{Interpretive cautions}
Stronger doublet member; stellar-wind contamination and unresolved saturation are common in galaxy spectra.

\subsection*{Useful companions}
Si IV 1402.770 A; C IV doublet; Si II 1260/1527 A

\subsection*{Sources}
\texttt{NIST\_ASD; MORTON\_2003; TUMLINSON\_2017}

\clearpage


\section{Si IV 1403}
\index{Si IV}

\begin{conceptbox}[At a glance]
\begin{tabularx}{\textwidth}{@{}lX@{}}
\textbf{Rest coordinate} & 1402.77 angstrom \\
\textbf{Band} & far-ultraviolet \\
\textbf{Transition} & 3s 2S1/2 -> 3p 2P1/2 \\
\textbf{Mechanism} & permitted resonance transition \\
\textbf{Primary tracers} & doublet saturation; ionized-gas kinematics; covering fraction \\
\textbf{Typical environments} & photoionized CGM; galactic winds; IGM absorbers; hot-star winds \\
\textbf{Indicative regime} & CIE peak around 6-8x10\textasciicircum{}4 K, or cooler photoionized gas \\
\textbf{Density guidance} & Not a direct density probe; doublet ratio and profile constrain opacity and covering. \\
\end{tabularx}
\end{conceptbox}

Si IV 1403 is produced by permitted resonance transition and is encountered in photoionized CGM; galactic winds; IGM absorbers; hot-star winds. It is commonly used to trace doublet saturation; ionized-gas kinematics; covering fraction.

\subsection*{How to use it}
Use this transition with its diagnostic companions to constrain doublet saturation; ionized-gas kinematics; covering fraction, while fitting the continuum, line-spread function, and relevant blends.

\subsection*{Interpretive cautions}
Weaker member helps reveal saturation; nearby interstellar/stellar features must be modeled in integrated spectra.

\subsection*{Useful companions}
Si IV 1393.755 A; C IV doublet; O VI doublet

\subsection*{Sources}
\texttt{NIST\_ASD; MORTON\_2003; TUMLINSON\_2017}

\clearpage


\section{C II 1334}
\index{C II}

\begin{conceptbox}[At a glance]
\begin{tabularx}{\textwidth}{@{}lX@{}}
\textbf{Rest coordinate} & 1334.5323 angstrom \\
\textbf{Band} & far-ultraviolet \\
\textbf{Transition} & 2s2 2p 2P1/2 -> 2s2 3s 2S1/2 \\
\textbf{Mechanism} & permitted resonance transition \\
\textbf{Primary tracers} & low-ionization gas column and kinematics; metallicity with ionization/depletion corrections; outflow covering fraction \\
\textbf{Typical environments} & neutral and weakly ionized ISM; cool CGM; damped Ly-alpha systems; galactic outflows \\
\textbf{Indicative regime} & typically cool/warm gas around 10\textasciicircum{}2-10\textasciicircum{}4 K \\
\textbf{Density guidance} & Ground-state absorption is column sensitive; excited C II* 1335 lines provide density/cooling constraints when separately measured. \\
\end{tabularx}
\end{conceptbox}

C II 1334 is produced by permitted resonance transition and is encountered in neutral and weakly ionized ISM; cool CGM; damped Ly-alpha systems; galactic outflows. It is commonly used to trace low-ionization gas column and kinematics; metallicity with ionization/depletion corrections; outflow covering fraction.

\subsection*{How to use it}
Use this transition with its diagnostic companions to constrain low-ionization gas column and kinematics; metallicity with ionization/depletion corrections; outflow covering fraction, while fitting the continuum, line-spread function, and relevant blends.

\subsection*{Interpretive cautions}
Often saturated even when not black because of unresolved velocity structure or partial covering; emission infill can reduce apparent absorption.

\subsection*{Useful companions}
C II* 1335.66/1335.71 A; Si II 1260 A; O I 1302 A; [C II] 158 micron

\subsection*{Sources}
\texttt{NIST\_ASD; MORTON\_2003; TUMLINSON\_2017}

\clearpage


\section{H I Ly-alpha}
\index{H I}

\begin{conceptbox}[At a glance]
\begin{tabularx}{\textwidth}{@{}lX@{}}
\textbf{Rest coordinate} & 1215.6701 angstrom \\
\textbf{Band} & far-ultraviolet \\
\textbf{Transition} & Lyman n=2 -> 1 (fine-structure blend) \\
\textbf{Mechanism} & permitted resonant transition; recombination/collisional emission and absorption \\
\textbf{Primary tracers} & neutral-hydrogen column; IGM density field; reionization topology statistically; gas kinematics through radiative-transfer modeling \\
\textbf{Typical environments} & H II regions and galaxies; CGM; Ly-alpha forest IGM; damped Ly-alpha systems; AGN \\
\textbf{Indicative regime} & neutral/photoionized gas around 10\textasciicircum{}4 K is common, but absorption spans broad thermal histories \\
\textbf{Density guidance} & Not a direct density line; forest statistics and damping wings constrain density/column after ionization and radiative-transfer modeling. \\
\end{tabularx}
\end{conceptbox}

H I Ly-alpha is produced by permitted resonant transition; recombination/collisional emission and absorption and is encountered in H II regions and galaxies; CGM; Ly-alpha forest IGM; damped Ly-alpha systems; AGN. It is commonly used to trace neutral-hydrogen column; IGM density field; reionization topology statistically; gas kinematics through radiative-transfer modeling.

\subsection*{How to use it}
Use this transition with its diagnostic companions to constrain neutral-hydrogen column; IGM density field; reionization topology statistically; gas kinematics through radiative-transfer modeling, while fitting the continuum, line-spread function, and relevant blends.

\subsection*{Interpretive cautions}
Resonant scattering, dust and IGM absorption decouple observed flux/profile from intrinsic production. Geocoronal Ly-alpha contaminates low-redshift work.

\subsection*{Useful companions}
Ly-beta 1025.72 A; H I Lyman limit; metal absorption lines; H-alpha under recombination assumptions

\subsection*{Sources}
\texttt{NIST\_ASD; DRAINE\_2011; TUMLINSON\_2017}

\clearpage


\section{O VI 1032}
\index{O VI}

\begin{conceptbox}[At a glance]
\begin{tabularx}{\textwidth}{@{}lX@{}}
\textbf{Rest coordinate} & 1031.9261 angstrom \\
\textbf{Band} & far-ultraviolet \\
\textbf{Transition} & 2s 2S1/2 -> 2p 2P3/2 \\
\textbf{Mechanism} & permitted resonance transition \\
\textbf{Primary tracers} & highly ionized oxygen; transition-temperature gas; feedback and halo baryons; outflow kinematics \\
\textbf{Typical environments} & warm-hot CGM/IGM; interfaces and mixing layers; radiative shocks; galactic winds; hot-star winds \\
\textbf{Indicative regime} & CIE abundance peaks near 3x10\textasciicircum{}5 K; photoionized O VI may occur at much lower density and temperature \\
\textbf{Density guidance} & Absorption gives column rather than density; combine with H I, path length, aligned ions and models. Emission scales strongly with density. \\
\end{tabularx}
\end{conceptbox}

O VI 1032 is produced by permitted resonance transition and is encountered in warm-hot CGM/IGM; interfaces and mixing layers; radiative shocks; galactic winds; hot-star winds. It is commonly used to trace highly ionized oxygen; transition-temperature gas; feedback and halo baryons; outflow kinematics.

\subsection*{How to use it}
Use this transition with its diagnostic companions to constrain highly ionized oxygen; transition-temperature gas; feedback and halo baryons; outflow kinematics, while fitting the continuum, line-spread function, and relevant blends.

\subsection*{Interpretive cautions}
Stronger member. H2 and interstellar blends, saturation, non-equilibrium cooling and multiphase alignment complicate temperature inference.

\subsection*{Useful companions}
O VI 1037.6167 A; H I Ly-beta nearby; C IV/Ne VIII; X-ray O VII

\subsection*{Sources}
\texttt{NIST\_ASD; MORTON\_2003; TUMLINSON\_2017}

\clearpage


\section{O VI 1038}
\index{O VI}

\begin{conceptbox}[At a glance]
\begin{tabularx}{\textwidth}{@{}lX@{}}
\textbf{Rest coordinate} & 1037.6167 angstrom \\
\textbf{Band} & far-ultraviolet \\
\textbf{Transition} & 2s 2S1/2 -> 2p 2P1/2 \\
\textbf{Mechanism} & permitted resonance transition \\
\textbf{Primary tracers} & O VI saturation/covering; transition-temperature gas; high-ionization kinematics \\
\textbf{Typical environments} & warm-hot CGM/IGM; interfaces; shocks; winds \\
\textbf{Indicative regime} & CIE peak near 3x10\textasciicircum{}5 K, with viable photoionized alternatives \\
\textbf{Density guidance} & No direct density measurement; doublet consistency reveals saturation and blends. \\
\end{tabularx}
\end{conceptbox}

O VI 1038 is produced by permitted resonance transition and is encountered in warm-hot CGM/IGM; interfaces; shocks; winds. It is commonly used to trace O VI saturation/covering; transition-temperature gas; high-ionization kinematics.

\subsection*{How to use it}
Use this transition with its diagnostic companions to constrain O VI saturation/covering; transition-temperature gas; high-ionization kinematics, while fitting the continuum, line-spread function, and relevant blends.

\subsection*{Interpretive cautions}
Weaker member is frequently blended with C II 1036, C II* and Galactic H2; fit all contaminants jointly.

\subsection*{Useful companions}
O VI 1031.9261 A; C II 1036.34 A; H2 absorption lines; Ne VIII doublet

\subsection*{Sources}
\texttt{NIST\_ASD; MORTON\_2003; TUMLINSON\_2017}

\clearpage


\section{Ne VIII 770}
\index{Ne VIII}

\begin{conceptbox}[At a glance]
\begin{tabularx}{\textwidth}{@{}lX@{}}
\textbf{Rest coordinate} & 770.409 angstrom \\
\textbf{Band} & extreme/far ultraviolet \\
\textbf{Transition} & 2s 2S1/2 -> 2p 2P3/2 \\
\textbf{Mechanism} & permitted resonance transition \\
\textbf{Primary tracers} & hotter ionized gas than O VI; warm-hot baryons; multiphase halo interfaces \\
\textbf{Typical environments} & warm-hot CGM/IGM; quasar absorbers; conductive/turbulent interfaces; AGN outflows \\
\textbf{Indicative regime} & CIE abundance peaks near 6-7x10\textasciicircum{}5 K; low-density photoionized solutions require long path lengths \\
\textbf{Density guidance} & Column/ionization diagnostic, not a direct density line; combine with O VI, H I and adjacent ions. \\
\end{tabularx}
\end{conceptbox}

Ne VIII 770 is produced by permitted resonance transition and is encountered in warm-hot CGM/IGM; quasar absorbers; conductive/turbulent interfaces; AGN outflows. It is commonly used to trace hotter ionized gas than O VI; warm-hot baryons; multiphase halo interfaces.

\subsection*{How to use it}
Use this transition with its diagnostic companions to constrain hotter ionized gas than O VI; warm-hot baryons; multiphase halo interfaces, while fitting the continuum, line-spread function, and relevant blends.

\subsection*{Interpretive cautions}
Usually accessible only after redshifting above the Lyman limit/instrument cutoff. Forest blending and continuum placement dominate weak detections.

\subsection*{Useful companions}
Ne VIII 780.324 A; O VI 1032/1038 A; Mg X 610/625 A; X-ray O VII

\subsection*{Sources}
\texttt{NIST\_ASD; MORTON\_2003; TUMLINSON\_2017}

\clearpage


\section{Ne VIII 780}
\index{Ne VIII}

\begin{conceptbox}[At a glance]
\begin{tabularx}{\textwidth}{@{}lX@{}}
\textbf{Rest coordinate} & 780.324 angstrom \\
\textbf{Band} & extreme/far ultraviolet \\
\textbf{Transition} & 2s 2S1/2 -> 2p 2P1/2 \\
\textbf{Mechanism} & permitted resonance transition \\
\textbf{Primary tracers} & doublet confirmation; warm-hot ionized gas; saturation \\
\textbf{Typical environments} & warm-hot CGM/IGM; quasar absorbers; interfaces; AGN outflows \\
\textbf{Indicative regime} & CIE peak near 6-7x10\textasciicircum{}5 K; photoionized alternatives possible \\
\textbf{Density guidance} & Not direct; weaker doublet member constrains saturation and false-positive blends. \\
\end{tabularx}
\end{conceptbox}

Ne VIII 780 is produced by permitted resonance transition and is encountered in warm-hot CGM/IGM; quasar absorbers; interfaces; AGN outflows. It is commonly used to trace doublet confirmation; warm-hot ionized gas; saturation.

\subsection*{How to use it}
Use this transition with its diagnostic companions to constrain doublet confirmation; warm-hot ionized gas; saturation, while fitting the continuum, line-spread function, and relevant blends.

\subsection*{Interpretive cautions}
Require velocity-coincident detection with the stronger member when possible; the far-UV forest yields a significant interloper rate.

\subsection*{Useful companions}
Ne VIII 770.409 A; O VI doublet; X-ray O VII/O VIII

\subsection*{Sources}
\texttt{NIST\_ASD; MORTON\_2003; TUMLINSON\_2017}

\clearpage


\section{O VII He-alpha resonance (w)}
\index{O VII}

\begin{conceptbox}[At a glance]
\begin{tabularx}{\textwidth}{@{}lX@{}}
\textbf{Rest coordinate} & 21.6019 angstrom \\
\textbf{Band} & soft X-ray \\
\textbf{Transition} & 1s2 1S0 -> 1s2p 1P1 resonance line \\
\textbf{Mechanism} & He-like resonance absorption/emission, usually collisional or photoionized plasma \\
\textbf{Primary tracers} & million-kelvin plasma; missing baryons; hot-halo column; temperature/ionization with O VIII \\
\textbf{Typical environments} & hot CGM/IGM; galaxy groups and clusters; stellar coronae; supernova remnants; AGN warm absorbers \\
\textbf{Indicative regime} & CIE abundance peaks around 2x10\textasciicircum{}6 K; photoionized O VII can occur at lower kinetic temperature \\
\textbf{Density guidance} & The resonance line alone is not a density meter; the forbidden/intercombination ratio in the full He-like triplet is density and UV-field sensitive. \\
\end{tabularx}
\end{conceptbox}

O VII He-alpha resonance (w) is produced by He-like resonance absorption/emission, usually collisional or photoionized plasma and is encountered in hot CGM/IGM; galaxy groups and clusters; stellar coronae; supernova remnants; AGN warm absorbers. It is commonly used to trace million-kelvin plasma; missing baryons; hot-halo column; temperature/ionization with O VIII.

\subsection*{How to use it}
Use this transition with its diagnostic companions to constrain million-kelvin plasma; missing baryons; hot-halo column; temperature/ionization with O VIII, while fitting the continuum, line-spread function, and relevant blends.

\subsection*{Interpretive cautions}
Weak absorption requires high spectral resolution and careful treatment of instrumental features, Galactic foreground, look-elsewhere effects and wavelength calibration.

\subsection*{Useful companions}
O VII intercombination near 21.80 A; O VII forbidden 22.10 A; O VIII 18.967 A; O VI UV doublet

\subsection*{Sources}
\texttt{ATOMDB; NIST\_ASD; TUMLINSON\_2017}

\clearpage


\section{O VIII Ly-alpha1}
\index{O VIII}

\begin{conceptbox}[At a glance]
\begin{tabularx}{\textwidth}{@{}lX@{}}
\textbf{Rest coordinate} & 18.96711 angstrom \\
\textbf{Band} & soft X-ray \\
\textbf{Transition} & 2p 2P3/2 -> 1s 2S1/2 \\
\textbf{Mechanism} & hydrogenic resonance transition \\
\textbf{Primary tracers} & several-million-kelvin gas; ionization temperature with O VII; hot-gas kinematics \\
\textbf{Typical environments} & hot CGM/IGM; galaxy groups and clusters; stellar coronae; supernova remnants; AGN warm absorbers \\
\textbf{Indicative regime} & CIE abundance peaks around 3x10\textasciicircum{}6 K; photoionized gas is also common \\
\textbf{Density guidance} & Primarily temperature/ionization and column sensitive; density requires metastable lines, emission measure or photoionization modeling. \\
\end{tabularx}
\end{conceptbox}

O VIII Ly-alpha1 is produced by hydrogenic resonance transition and is encountered in hot CGM/IGM; galaxy groups and clusters; stellar coronae; supernova remnants; AGN warm absorbers. It is commonly used to trace several-million-kelvin gas; ionization temperature with O VII; hot-gas kinematics.

\subsection*{How to use it}
Use this transition with its diagnostic companions to constrain several-million-kelvin gas; ionization temperature with O VII; hot-gas kinematics, while fitting the continuum, line-spread function, and relevant blends.

\subsection*{Interpretive cautions}
The Ly-alpha doublet is often unresolved. Resonant scattering can suppress central emission in optically thick hot atmospheres.

\subsection*{Useful companions}
O VIII Ly-alpha2 18.9725 A; O VII triplet near 21.6-22.1 A; O VIII Ly-beta

\subsection*{Sources}
\texttt{ATOMDB; NIST\_ASD}

\clearpage


\section{Fe XVII 15.01}
\index{Fe XVII}

\begin{conceptbox}[At a glance]
\begin{tabularx}{\textwidth}{@{}lX@{}}
\textbf{Rest coordinate} & 15.014 angstrom \\
\textbf{Band} & soft X-ray \\
\textbf{Transition} & 2p6 1S0 -> 2p5 3d 1P1 (3C resonance line) \\
\textbf{Mechanism} & collisionally excited resonance emission \\
\textbf{Primary tracers} & roughly 5 MK plasma; resonant scattering/turbulence using Fe XVII ratios; cooling gas \\
\textbf{Typical environments} & stellar coronae; elliptical-galaxy atmospheres; cool-core clusters; supernova remnants \\
\textbf{Indicative regime} & CIE emissivity is strong around 4-7x10\textasciicircum{}6 K \\
\textbf{Density guidance} & The 15.01/17.05-17.10 complex is more sensitive to resonant scattering and atomic physics than ordinary low-density plasma density; some Fe XVII ratios gain density sensitivity at high densities. \\
\end{tabularx}
\end{conceptbox}

Fe XVII 15.01 is produced by collisionally excited resonance emission and is encountered in stellar coronae; elliptical-galaxy atmospheres; cool-core clusters; supernova remnants. It is commonly used to trace roughly 5 MK plasma; resonant scattering/turbulence using Fe XVII ratios; cooling gas.

\subsection*{How to use it}
Use this transition with its diagnostic companions to constrain roughly 5 MK plasma; resonant scattering/turbulence using Fe XVII ratios; cooling gas, while fitting the continuum, line-spread function, and relevant blends.

\subsection*{Interpretive cautions}
Resonant scattering lowers the central 15.01-A intensity and can constrain turbulence, but plasma-code and line-blend uncertainties must be propagated.

\subsection*{Useful companions}
Fe XVII 15.26 A; Fe XVII 16.78/17.05/17.10 A; O VIII 18.97 A

\subsection*{Sources}
\texttt{ATOMDB; NIST\_ASD}

\clearpage


\section{Ne IX He-alpha resonance (w)}
\index{Ne IX}

\begin{conceptbox}[At a glance]
\begin{tabularx}{\textwidth}{@{}lX@{}}
\textbf{Rest coordinate} & 13.4473 angstrom \\
\textbf{Band} & soft X-ray \\
\textbf{Transition} & 1s2 1S0 -> 1s2p 1P1 \\
\textbf{Mechanism} & He-like resonance transition \\
\textbf{Primary tracers} & few-million-kelvin plasma; ionization state with Ne X; plasma excitation with triplet ratios \\
\textbf{Typical environments} & stellar coronae; supernova remnants; hot galactic gas; AGN photoionized plasma \\
\textbf{Indicative regime} & CIE abundance/emissivity is strong around 4x10\textasciicircum{}6 K \\
\textbf{Density guidance} & Use the forbidden/intercombination components of the Ne IX triplet for density/UV-field sensitivity; w alone is not sufficient. \\
\end{tabularx}
\end{conceptbox}

Ne IX He-alpha resonance (w) is produced by He-like resonance transition and is encountered in stellar coronae; supernova remnants; hot galactic gas; AGN photoionized plasma. It is commonly used to trace few-million-kelvin plasma; ionization state with Ne X; plasma excitation with triplet ratios.

\subsection*{How to use it}
Use this transition with its diagnostic companions to constrain few-million-kelvin plasma; ionization state with Ne X; plasma excitation with triplet ratios, while fitting the continuum, line-spread function, and relevant blends.

\subsection*{Interpretive cautions}
Fe XIX-XXI blends can contaminate the triplet in hot plasmas; use a full plasma model at finite spectral resolution.

\subsection*{Useful companions}
Ne IX intercombination near 13.55 A; Ne IX forbidden 13.70 A; Ne X Ly-alpha 12.13 A

\subsection*{Sources}
\texttt{ATOMDB; NIST\_ASD}

\clearpage


\section{Fe XXV He-alpha resonance (w)}
\index{Fe XXV}

\begin{conceptbox}[At a glance]
\begin{tabularx}{\textwidth}{@{}lX@{}}
\textbf{Rest coordinate} & 1.8504 angstrom \\
\textbf{Band} & hard X-ray / Fe K complex \\
\textbf{Transition} & 1s2 1S0 -> 1s2p 1P1 \\
\textbf{Mechanism} & He-like resonance transition in highly ionized plasma \\
\textbf{Primary tracers} & very hot thermal plasma; bulk/turbulent velocity; ionization temperature; resonant scattering \\
\textbf{Typical environments} & galaxy clusters; accreting compact objects; solar/stellar flares; supernova remnants; AGN winds \\
\textbf{Indicative regime} & CIE line emissivity is strong at several 10\textasciicircum{}7 K \\
\textbf{Density guidance} & The resolved He-like complex and satellite ratios diagnose temperature, ionization age and at high density density; CCD-blended centroids are model dependent. \\
\end{tabularx}
\end{conceptbox}

Fe XXV He-alpha resonance (w) is produced by He-like resonance transition in highly ionized plasma and is encountered in galaxy clusters; accreting compact objects; solar/stellar flares; supernova remnants; AGN winds. It is commonly used to trace very hot thermal plasma; bulk/turbulent velocity; ionization temperature; resonant scattering.

\subsection*{How to use it}
Use this transition with its diagnostic companions to constrain very hot thermal plasma; bulk/turbulent velocity; ionization temperature; resonant scattering, while fitting the continuum, line-spread function, and relevant blends.

\subsection*{Interpretive cautions}
Quote whether energy refers to the w line or an unresolved Fe XXV complex. Gain calibration and plasma-code choices set velocity/temperature systematics.

\subsection*{Useful companions}
Fe XXV x,y,z components; dielectronic satellite lines; Fe XXVI Ly-alpha; neutral Fe K-alpha

\subsection*{Sources}
\texttt{ATOMDB; NIST\_ASD}

\clearpage


\section{Fe XXVI Ly-alpha1}
\index{Fe XXVI}

\begin{conceptbox}[At a glance]
\begin{tabularx}{\textwidth}{@{}lX@{}}
\textbf{Rest coordinate} & 1.778 angstrom \\
\textbf{Band} & hard X-ray / Fe K complex \\
\textbf{Transition} & 2p 2P3/2 -> 1s 2S1/2 \\
\textbf{Mechanism} & hydrogenic resonance transition \\
\textbf{Primary tracers} & extreme ionization; very hot plasma; ultrafast outflow velocity; ionization temperature with Fe XXV \\
\textbf{Typical environments} & hottest cluster plasma; accreting black-hole and neutron-star winds; solar/stellar flares; photoionized AGN \\
\textbf{Indicative regime} & CIE fraction becomes important around 10\textasciicircum{}8 K; photoionized Fe XXVI also occurs \\
\textbf{Density guidance} & Primarily ionization/temperature and column sensitive; density comes from broader photoionization or time-dependent modeling. \\
\end{tabularx}
\end{conceptbox}

Fe XXVI Ly-alpha1 is produced by hydrogenic resonance transition and is encountered in hottest cluster plasma; accreting black-hole and neutron-star winds; solar/stellar flares; photoionized AGN. It is commonly used to trace extreme ionization; very hot plasma; ultrafast outflow velocity; ionization temperature with Fe XXV.

\subsection*{How to use it}
Use this transition with its diagnostic companions to constrain extreme ionization; very hot plasma; ultrafast outflow velocity; ionization temperature with Fe XXV, while fitting the continuum, line-spread function, and relevant blends.

\subsection*{Interpretive cautions}
Often unresolved as a doublet; line identification in fast AGN winds can be degenerate with Fe XXV at a different velocity.

\subsection*{Useful companions}
Fe XXVI Ly-alpha2 1.7835 A; Fe XXV 1.8504 A; Fe XXVI Ly-beta

\subsection*{Sources}
\texttt{ATOMDB; NIST\_ASD}

\clearpage


\section{neutral Fe K-alpha1}
\index{Fe I (and low-ionization Fe)}

\begin{conceptbox}[At a glance]
\begin{tabularx}{\textwidth}{@{}lX@{}}
\textbf{Rest coordinate} & 6.4040062 keV \\
\textbf{Band} & hard X-ray \\
\textbf{Transition} & K-L3 fluorescence, 2p3/2 -> 1s vacancy \\
\textbf{Mechanism} & K-shell fluorescence following photoionization or energetic-particle impact \\
\textbf{Primary tracers} & cold reflection; reprocessing geometry; black-hole spin through relativistic broadening with a disk model; past X-ray illumination \\
\textbf{Typical environments} & AGN tori and accretion disks; X-ray binaries; molecular clouds irradiated by X-rays; stellar flares \\
\textbf{Indicative regime} & cold or weakly ionized material illuminated by photons above the Fe K edge; not a thermal-excitation temperature line \\
\textbf{Density guidance} & Equivalent width and Compton shoulder constrain column, abundance and geometry rather than local density uniquely. \\
\end{tabularx}
\end{conceptbox}

neutral Fe K-alpha1 is produced by K-shell fluorescence following photoionization or energetic-particle impact and is encountered in AGN tori and accretion disks; X-ray binaries; molecular clouds irradiated by X-rays; stellar flares. It is commonly used to trace cold reflection; reprocessing geometry; black-hole spin through relativistic broadening with a disk model; past X-ray illumination.

\subsection*{How to use it}
Use this transition with its diagnostic companions to constrain cold reflection; reprocessing geometry; black-hole spin through relativistic broadening with a disk model; past X-ray illumination, while fitting the continuum, line-spread function, and relevant blends.

\subsection*{Interpretive cautions}
The common '6.4 keV line' is a K-alpha doublet and shifts with ionization. Broad profiles require continuum, absorption and relativistic-reflection systematics to be modeled jointly.

\subsection*{Useful companions}
Fe K-alpha2 6.391026 keV; Fe K-beta near 7.06 keV; Fe K edge; Compton reflection hump; Fe XXV/XXVI

\subsection*{Sources}
\texttt{NIST\_FE\_K; NIST\_XRAY}

\clearpage


\section{electron-positron annihilation line}
\index{e- + e+}

\begin{conceptbox}[At a glance]
\begin{tabularx}{\textwidth}{@{}lX@{}}
\textbf{Rest coordinate} & 511 keV \\
\textbf{Band} & gamma ray \\
\textbf{Transition} & two-photon annihilation at the electron rest-mass energy \\
\textbf{Mechanism} & direct annihilation and para-positronium two-photon decay \\
\textbf{Primary tracers} & positron production and transport; annihilation phase of the ISM; compact-source activity statistically \\
\textbf{Typical environments} & Galactic bulge and disk ISM; compact-object environments; radioactive-ejecta environments \\
\textbf{Indicative regime} & line width and positronium fraction probe warm neutral/ionized annihilation media; production is not tied to gas temperature \\
\textbf{Density guidance} & No critical density; annihilation timescale and positronium fraction depend on phase, ionization fraction, grains and density. \\
\end{tabularx}
\end{conceptbox}

electron-positron annihilation line is produced by direct annihilation and para-positronium two-photon decay and is encountered in Galactic bulge and disk ISM; compact-object environments; radioactive-ejecta environments. It is commonly used to trace positron production and transport; annihilation phase of the ISM; compact-source activity statistically.

\subsection*{How to use it}
Use this transition with its diagnostic companions to constrain positron production and transport; annihilation phase of the ISM; compact-source activity statistically, while fitting the continuum, line-spread function, and relevant blends.

\subsection*{Interpretive cautions}
The centroid is fixed by rest mass, while astrophysical broadening and the positronium continuum constrain thermalization/annihilation conditions. Diffuse-background modeling is central.

\subsection*{Useful companions}
three-photon positronium continuum below 511 keV; 26Al 1808.65 keV; 60Fe 1173/1333 keV

\subsection*{Sources}
\texttt{NASA\_511}

\clearpage


\section{26Al decay gamma-ray line}
\index{26Mg* following 26Al decay}

\begin{conceptbox}[At a glance]
\begin{tabularx}{\textwidth}{@{}lX@{}}
\textbf{Rest coordinate} & 1808.65 keV \\
\textbf{Band} & gamma ray \\
\textbf{Transition} & de-excitation of the 1808.65-keV state of 26Mg \\
\textbf{Mechanism} & radioactive beta-plus/electron-capture decay followed by nuclear gamma de-excitation \\
\textbf{Primary tracers} & ongoing nucleosynthesis; recent massive-star activity over Myr timescales; hot-ejecta kinematics and Galactic rotation \\
\textbf{Typical environments} & Galactic plane; massive-star associations; superbubbles; core-collapse supernova and Wolf-Rayet ejecta \\
\textbf{Indicative regime} & decay is independent of thermal excitation; emitting isotope is transported through hot superbubbles and ISM phases \\
\textbf{Density guidance} & No critical density; intensity traces isotope mass and decay, while centroid/width trace bulk motion and turbulence. \\
\end{tabularx}
\end{conceptbox}

26Al decay gamma-ray line is produced by radioactive beta-plus/electron-capture decay followed by nuclear gamma de-excitation and is encountered in Galactic plane; massive-star associations; superbubbles; core-collapse supernova and Wolf-Rayet ejecta. It is commonly used to trace ongoing nucleosynthesis; recent massive-star activity over Myr timescales; hot-ejecta kinematics and Galactic rotation.

\subsection*{How to use it}
Use this transition with its diagnostic companions to constrain ongoing nucleosynthesis; recent massive-star activity over Myr timescales; hot-ejecta kinematics and Galactic rotation, while fitting the continuum, line-spread function, and relevant blends.

\subsection*{Interpretive cautions}
Laboratory energy is 1808.65 +/- 0.07 keV. Instrumental background, broad response and diffuse-sky modeling dominate; the mean lifetime is about 1.04 Myr.

\subsection*{Useful companions}
60Fe 1173 and 1333 keV; 511 keV annihilation; massive-star tracers and free-free emission

\subsection*{Sources}
\texttt{SAMWORTH\_1972; DIEHL\_2021}

\clearpage
